% ════════════════════════════════════════════════════════════════
% Zestaw przykładowych instancji do powtórki przed egzaminem.
% Same DANE wejściowe (po trzy instancje na algorytm) --- bez rozwiązań,
% do samodzielnego prześledzenia działania algorytmu.
% Kompilacja: lualatex (patrz ../preamble, build.sh).
% ════════════════════════════════════════════════════════════════

\documentclass[11pt,a4paper]{article}

%\usepackage[utf8]{inputenc}
%\usepackage[T1]{fontenc}
\usepackage[polish]{babel}
\usepackage{geometry}
\usepackage{booktabs}
\usepackage{tabularray}
\usepackage{afterpage}
\usepackage{rotating}
\usepackage{longtable}
\usepackage{array}
\usepackage{xcolor}
\usepackage{hyperref}
\usepackage{titlesec}
\usepackage{parskip}
\usepackage{enumitem}
\usepackage{fancyhdr}
\usepackage{pdflscape}
\usepackage{makecell}
\usepackage[normalem]{ulem}% normalem: nie przejmuj \emph; \uline łamie się między wierszami

\usepackage{tikz}
\usetikzlibrary{decorations}
\usetikzlibrary{decorations.pathreplacing}
\usetikzlibrary{patterns}

\usepackage{tcolorbox}
\tcbuselibrary{breakable,skins}

\usepackage{caption}

\usepackage{mathtools}
\usepackage{etoolbox}
\usepackage{xpatch}
\usepackage{environ}% \RenewEnviron: przechwycenie treści dowodu do pominięcia (AZ_SKIP)

\usepackage[noend]{algpseudocode}
\usepackage{algorithm}
\usepackage{multirow}
\floatname{algorithm}{Algorytm}
\renewcommand{\listalgorithmname}{Spis algorytmów}

\usepackage{nicematrix}
\usepackage{cancel}
\usepackage{amsthm}
%\usepackage{amssymb}
\usepackage{amsmath}
\usepackage[warnings-off={mathtools-overbracket,mathtools-colon}]{unicode-math}
\setmathfont{XITS Math}
\mathtoolsset{mathic=true}

% Wariant ⩽/⩾ z linią równości pod tym samym kątem co < / > (jak \leqslant/\geqslant).
% Remap odłożony do \begin{document}, bo unicode-math finalizuje tablicę symboli
% dopiero przy starcie dokumentu i nadpisałby zwykły \let z preambuły.
\AtBeginDocument{%
	\let\leq\leqslant
	\let\geq\geqslant
	\allowdisplaybreaks
}

\geometry{margin=2.5cm, headheight=12.7pt}

% Ostatnia deska ratunku przy łamaniu akapitów: pozwala TeX-owi rozciągnąć
% trudny wiersz zamiast wypychać go poza margines (usuwa drobne Overfull \hbox
% w gęstym tekście z wtrąceniami matematycznymi).
\emergencystretch=2em
\hfuzz=0.6pt% toleruj przepełnienia poniżej ~0,2 mm (niewidoczne, a generują ostrzeżenia)

\definecolor{accent}{RGB}{30, 60, 100}
\definecolor{accentA}{RGB}{255, 187, 0}
\definecolor{accentB}{RGB}{255, 100, 0}
\definecolor{accentC}{RGB}{0, 140, 130}
\definecolor{placeholdercolor}{RGB}{160, 160, 160}
\definecolor{rowA}{RGB}{235, 242, 250}
\definecolor{dygresjagray}{RGB}{105, 105, 105}% treść dygresji — wyszarzona, materiał opcjonalny
\definecolor{obowiazkowybg}{RGB}{255, 252, 245}% tło dowodów obowiązkowych na egzaminie
\definecolor{obowiazkowyline}{RGB}{230, 200, 150}% nasycona linia na granicy podświetlenia
\definecolor{pomijalnybg}{RGB}{250, 250, 250}% tło dowodów pomijalnych przy powtórce
\definecolor{pomijalnyline}{RGB}{240, 240, 240}% szara linia na granicy podświetlenia
\definecolor{pomijalnytext}{RGB}{120, 120, 120}% przygaszony tekst dowodów pomijalnych

% Podświetlenie materiału obowiązkowego jako zamknięte pudełko w bloku tekstu
% (nie rozlewa się na marginesy). Delikatna ramka + akcentowy pasek z lewej dają
% czytelny sygnał „to obowiązkowe” bez agresywnego tła do krawędzi kartki.
% breakable: poprawnie łamie się między stronami (enhanced rysuje ładne górne/
% dolne krawędzie segmentów). Stranded rysunek może zostawić trochę tła u dołu
% łamanego segmentu, ale w zamkniętym pudełku jest to znacznie mniej rażące.
\newtcolorbox{obowiazkowybox}{%
	breakable, enhanced,
	colback=obowiazkowybg,
	colframe=obowiazkowyline,
	boxrule=0.5pt, % jednolita cienka ramka ze wszystkich stron
	arc=2pt, outer arc=2pt,
	left=10pt, right=10pt, top=8pt, bottom=8pt,
	boxsep=0pt,
	pad at break*=6pt,
}
% Usuń resztkowy odstęp pionowy na końcu ramki (np. \topsep zamykający proof),
% który potrafi przelać się na kolejną stronę jako pusty segment tcolorbox
% i zostawić blady pasek tła (colback rysuje się dla każdego segmentu).
\AtEndEnvironment{obowiazkowybox}{\par\removelastskip}
% Zarezerwuj miejsce PRZED otwarciem ramki. Inaczej \azNeedSpace twierdzenia
% wewnątrz świeżo otwartej ramki wymusza łamanie strony przy pustej zawartości,
% tworząc pusty pierwszy segment (blady pasek tła u góry/dołu strony).
\BeforeBeginEnvironment{obowiazkowybox}{\Needspace*{6\baselineskip}}

% ── WERSJA Z WYRÓŻNIENIAMI ───────────────────────────────────────
% Flaga \azhighlight włącza tło materiału obowiązkowego. Sterowana zmienną
% środowiskową AZ_HIGHLIGHT=1 (czytaną przez LuaTeX przy starcie) — ten sam
% plik źródłowy daje wersję zwykłą i wersję z wyróżnieniami.
\newbool{azhighlight}
\edef\azHighlightEnv{\directlua{tex.sprint(os.getenv("AZ_HIGHLIGHT") or "0")}}
\ifdefstring{\azHighlightEnv}{1}{\booltrue{azhighlight}}{}

% Flaga \azskip (AZ_SKIP=1): wycina nieobowiązkowe dowody w całości — wersja do
% szybkiej powtórki, w której zostają tylko dowody obowiązkowe.
\newbool{azskip}
\edef\azSkipEnv{\directlua{tex.sprint(os.getenv("AZ_SKIP") or "0")}}
\ifdefstring{\azSkipEnv}{1}{\booltrue{azskip}}{}

% Środowisko obowiazkowy: w wersji z wyróżnieniami otwiera ramkę
% obowiazkowybox, w zwykłej wersji jest przezroczyste (treść składana
% normalnie, bez tła). Dzięki temu twierdzenia obowiązkowe oznaczamy raz,
% a obie wersje generujemy z jednego źródła.
\newenvironment{obowiazkowy}
{\ifbool{azhighlight}{\begin{obowiazkowybox}}{}}
			{\ifbool{azhighlight}{\end{obowiazkowybox}}{}}

% Dowody nieobowiązkowe: w wersji z wyróżnieniami szare tło sygnalizuje, że
% można je pominąć przy powtórce do egzaminu. Zamiast oznaczać każdy dowód
% ręcznie, w wersji z wyróżnieniami opakowujemy CAŁE środowisko proof w szare
% pudełko — z wyjątkiem dowodów leżących już w ramce obowiazkowy (te są
% istotne i mają tło cream). Flaga \c@inobow śledzi, czy jesteśmy w obowiazkowy.
% W wersji zwykłej \tcolorboxenvironment nie jest wołane → proof bez zmian.
% inobow śledzi „jesteśmy w materiale obowiązkowym”. Ustawiamy go na środowisku
% obowiazkowy (a nie obowiazkowybox), żeby działał też bez wyróżnień — wersja
% z pominięciem dowodów potrzebuje tej informacji niezależnie od AZ_HIGHLIGHT.
\newbool{inobow}
\AtBeginEnvironment{obowiazkowy}{\booltrue{inobow}}
\AtEndEnvironment{obowiazkowy}{\boolfalse{inobow}}
\tcbset{
	pomijalnystyle/.style={%
		breakable, enhanced,
		colback=pomijalnybg,
		colframe=pomijalnyline,
		coltext=pomijalnytext,
		boxrule=0.5pt,
		arc=2pt, outer arc=2pt,
		left=10pt, right=10pt, top=8pt, bottom=8pt,
		boxsep=0pt,
		pad at break*=6pt,
		before skip=10pt plus 2pt, after skip=10pt plus 2pt,
	},
	% dowód w ramce obowiazkowy: bez szarego pudełka (zostaje na tle cream)
	pomijalnypusty/.style={blanker, before skip=\topsep, after skip=\topsep},
	% Rozstrzygamy styl jako KOD (a nie listę kluczy) w chwili czytania opcji
	% pudełka — wtedy \ifbool widzi już ustawione inobow (otwarte obowiazkowybox),
	% a problem dzielenia listy po przecinkach znika.
	proofbox/.code={\ifbool{inobow}{\tcbset{pomijalnypusty}}{\tcbset{pomijalnystyle}}},
}
% AZ_SKIP wygrywa: nieobowiązkowe dowody znikają zupełnie (przechwytujemy treść
% przez \BODY i nie składamy jej). Obowiązkowe składamy normalnie oryginalnym
% środowiskiem proof. Inaczej — w trybie wyróżnień — szare pudełka jak wyżej.
\ifbool{azskip}{%
	\let\azoldproof\proof
	\let\azoldendproof\endproof
	\RenewEnviron{proof}[1][\proofname]{%
		\ifbool{inobow}{\azoldproof[#1]\BODY\azoldendproof}{}%
	}%
}{\ifbool{azhighlight}{\tcolorboxenvironment{proof}{proofbox}}{}}

\titleformat{\section}{\large\bfseries\color{accent}}{}{0em}{}[\titlerule]
\titleformat{\subsection}{\normalsize\bfseries}{}{0em}{}
\titleformat{\subsubsection}{\normalsize\itshape}{}{0em}{}

\pagestyle{fancy}
\fancyhf{}

\rhead{\small \rightmark}
\lhead{\small \textit{Algorytmy Zaawansowane}}
\cfoot{\thepage}
\renewcommand{\headrulewidth}{0.4pt}

\hypersetup{colorlinks=true, linkcolor=accent, urlcolor=accent}

% ── KOMENDY ──────────────────────────────────────────────────────
% Wspólny nagłówek twierdzeń: gdy podano nazwę (argument [opcjonalny]),
% wypisz ją w nawiasie i przełam wiersz przed treścią; bez nazwy nagłówek
% pozostaje w jednej linii (run-in). \azThmHead{name}{number}{note}.
\newcommand{\azThmHead}[3]{%
	\thmname{#1}\thmnumber{ #2}%
	\def\azNote{#3}%
	\ifx\azNote\empty. \else\ (#3).\newline\fi}

% Większy odstęp przed/po + wspólny nagłówek z przełamaniem dla nazwanych.
% Wariant „definition” (treść prosta) i „plain” (treść kursywą).
\newtheoremstyle{azdefinition}%
{14pt plus 3pt minus 2pt}{12pt plus 3pt minus 2pt}%  odstęp przed / po
{\normalfont}%                                        czcionka treści
{0pt}%                                                wcięcie
{\bfseries}%                                          czcionka nagłówka
{}{0pt}%                                              interpunkcja i odstęp obsłużone w \azThmHead (headsep = 0)
{\azThmHead{#1}{#2}{#3}}%                          nagłówek
\newtheoremstyle{azplain}%
{14pt plus 3pt minus 2pt}{12pt plus 3pt minus 2pt}%
{\itshape}%
{0pt}%
{\bfseries}%
{}{0pt}%
{\azThmHead{#1}{#2}{#3}}%

\theoremstyle{azdefinition}
\newtheorem{definition}{Definicja}[section]
\newtheorem{example}[definition]{Przykład}
\newtheorem{problem}[definition]{Problem}

\theoremstyle{azplain}
\newtheorem{properties}[definition]{Własności}
\newtheorem{theorem}[definition]{Twierdzenie}
\newtheorem{conclusion}[definition]{Wniosek}
\newtheorem{lemma}[definition]{Lemat}

\theoremstyle{remark}
\newtheorem*{remark}{Uwaga}

% Nie pozostawiaj osamotnionego nagłówka twierdzenia na końcu strony:
% zarezerwuj kilka wierszy, w razie potrzeby przełam stronę przed środowiskiem.
\usepackage{needspace}
\newcommand{\azNeedSpace}{\Needspace*{4\baselineskip}}
\AtBeginEnvironment{definition}{\azNeedSpace}
\AtBeginEnvironment{example}{\azNeedSpace}
\AtBeginEnvironment{problem}{\azNeedSpace}
\AtBeginEnvironment{properties}{\azNeedSpace}
\AtBeginEnvironment{theorem}{\azNeedSpace}
\AtBeginEnvironment{conclusion}{\azNeedSpace}
\AtBeginEnvironment{lemma}{\azNeedSpace}

% Dygresja: wolniejsze rozwinięcie kroku dla czytelnika, który gubi się w rachunkach.
% Wyszarzona treść sygnalizuje materiał opcjonalny (można pominąć przy powtórce).
\newtheoremstyle{dygresjastyle}%
{\topsep}{\topsep}%                       odstęp przed / po
{\color{dygresjagray}}%                    czcionka treści (wyszarzona)
{0pt}%                                     wcięcie
{\color{dygresjagray}\itshape}%            czcionka nagłówka
{.}{.5em}{}%                               interpunkcja, odstęp po nagłówku, spec.
\theoremstyle{dygresjastyle}
\newtheorem*{dygresja}{Dygresja}


% ── KOMENDY ──────────────────────────────────────────────────────
% Podkreślenie terminów, które łamie się między wierszami (zwykłe \underline nie łamie).
\newcommand{\uemph}[1]{\uline{#1}}

\newcommand{\N}[0]{\mathbb{N}}
\newcommand{\Np}[0]{\mathbb{N_+}}
\newcommand{\R}[0]{\mathbb{R}}
\newcommand{\Rp}[0]{\mathbb{R_+}}

\newcommand{\BO}[0]{\mathcal{O}}
\newcommand{\BM}[0]{\mathcal{\Omega}}
\newcommand{\BT}[0]{\mathcal{\Theta}}

% Podwójne przekreślenie w „zwykłym” kierunku (//) — pakiet cancel daje tylko
% \cancel (/), \bcancel (\) i \xcancel (X). Oznacza skrócenie z wyrazem
% schowanym w wielokropku (drugi partner nie jest pokazany).
\newcommand{\dcancel}[1]{{\mathrlap{\cancel{#1}}\mkern8mu\cancel{\phantom{#1}}}}
% To samo w kierunku wstecznym (\\) — dla kompletu.
\newcommand{\dbcancel}[1]{{\mathrlap{\bcancel{#1}}\mkern8mu\bcancel{\phantom{#1}}}}

% \binom z XITS Math nie skaluje nawiasów wbudowanych w \genfrac — dolny
% argument wystaje pod nawias. Owijamy bezkreskowy \genfrac w prawdziwe
% \left(...\right), których wysokość dopasowuje się do dwuwierszowej zawartości.
\renewcommand{\binom}[2]{\left(\genfrac{}{}{0pt}{}{#1}{#2}\right)}

% ── KOMENDY ──────────────────────────────────────────────────────
% change default thickness of brackets to .6pt
\MHInternalSyntaxOn
\xpatchcmd\upbracketfill
{\sbox\z@{$\braceld$}\edef\l_MT_bracketheight_fdim{\the\ht\z@}}
{\edef\l_MT_bracketheight_fdim{.6pt}}
{}{\fail}

\xpatchcmd\downbracketfill
{\sbox\z@{$\braceld$}\edef\l_MT_bracketheight_fdim{\the\ht\z@}}
{\edef\l_MT_bracketheight_fdim{.6pt}}
{}{\fail}
\MHInternalSyntaxOff

% ── KOMENDY ──────────────────────────────────────────────────────

\algnewcommand{\To}[0]{\textbf{to} }
\algnewcommand{\Assign}[0]{$\;\,\gets\,\;$}
\algnewcommand{\Length}[1]{\textit{length}(#1)}
\algnewcommand{\Swap}[2]{\textit{swap}(#1, #2)}
\algnewcommand{\Size}[1]{\textit{size}(#1)}

\makeatletter
\newcommand{\FunctionNN}[2]{%
	\let\OldAlgLineNum\alglinenumber%
	\renewcommand{\alglinenumber}[1]{}%
	\Function{#1}{#2}%
	\let\alglinenumber\OldAlgLineNum%
	\addtocounter{ALG@line}{-1}% 
}
\makeatother

\makeatletter
\newenvironment{alg}[3][]{% #1=optional tcolorbox options, #2=caption, #3=label
	\if@noskipsec \leavevmode \fi% flush a pending run-in heading (\paragraph) so it stays above the box
	\par
	\addvspace{\bigskipamount}% breathing room before the box
	\refstepcounter{algorithm}\label{#3}% step the float counter + set ref/hyperref target
	\addcontentsline{\csname ext@algorithm\endcsname}{algorithm}% list-of-algorithms entry
	{\protect\numberline{\thealgorithm}{\ignorespaces #2}}%
	\small%
	\begin{tcolorbox}[
			breakable,
			sharp corners,
			boxrule=1pt,
			colback=white,
			colframe=black,
			colbacktitle=accent,
			coltitle=white,
			fonttitle=\bfseries,
			title={\csname fnum@algorithm\endcsname:\ \mdseries #2},% "Algorytm N: caption"
			left=2mm,
			right=2mm,
			top=2mm,
			bottom=2mm,
			#1%                    pass through any caller-supplied options
		]%
		\begin{algorithmic}[1]%
			}{%
		\end{algorithmic}%
	\end{tcolorbox}%
	\par\addvspace{\bigskipamount}% breathing room after the box
}
\makeatother

% ── RYSUNKI I TABELE BEZ PŁYWANIA ────────────────────────────────
% Zamiast pływających środowisk figure/table (i tak używanych prawie zawsze
% z [H], co przekreśla sens pływania) — wstawiamy treść dokładnie tam, gdzie
% jest napisana. Dzięki temu rysunek/tabela może też znaleźć się WEWNĄTRZ
% wyróżnienia obowiazkowy (pływak by z niego uciekł). \captionsetup{type=...}
% sprawia, że zwykłe \caption / \caption* numeruje się i etykietuje (\label,
% spis) tak samo jak w pływaku — tyle że bez pływania.
\newenvironment{azfigure}{%
	\par\addvspace{\bigskipamount}%
	\begingroup
	\captionsetup{type=figure}%
	\centering
}{%
	\par
	\endgroup
	\addvspace{\bigskipamount}%
}
\newenvironment{aztable}{%
	\par\addvspace{\bigskipamount}%
	\begingroup
	\captionsetup{type=table}%
	\centering
}{%
	\par
	\endgroup
	\addvspace{\bigskipamount}%
}
% ─────────────────────────────────────────────────────────────────


\begin{document}

\begin{center}
	{\LARGE\bfseries\color{accent} Algorytmy Zaawansowane}\\[0.3cm]
	{\large\bfseries Przykładowe instancje do powtórki}\\[0.2cm]
	{\small\itshape Po trzy instancje na algorytm --- same dane wejściowe, do samodzielnego prześledzenia.}
\end{center}
\bigskip

\tableofcontents
\newpage

% ════════════════════════════════════════════════════════════════
\section{Para najbliższych punktów}
% ════════════════════════════════════════════════════════════════

\begin{example}[Najbliższe punkty -- instancja 1]
	$Q = \{(2,3),\ (12,30),\ (40,50),\ (5,1),\ (12,10),\ (3,4),\ (8,9),\ (1,1)\}$.\\
	Znaleźć parę punktów w~najmniejszej odległości (euklidesowej).
\end{example}

\begin{example}[Najbliższe punkty -- instancja 2]
	Znaleźć parę punktów w~najmniejszej odległości:
	\begin{aztable}
		\centering
		\renewcommand{\arraystretch}{1.2}
		\begin{NiceTabular}{c|cccccccc}[margin,hvlines]
			$p_i$ & 1 & 2 & 3 & 4 & 5 & 6 & 7 & 8 \\
			\hline
			$x$   & 0 & 1 & 4 & 5 & 7 & 8 & 9 & 11 \\
			$y$   & 0 & 6 & 2 & 5 & 1 & 7 & 3 & 6
		\end{NiceTabular}
	\end{aztable}
\end{example}

	\begin{example}[Najbliższe punkty -- instancja 3]
		$Q = \{(1,2),\ (4,4),\ (7,1),\ (8,8),\ (2,9),\ (10,3),\ (6,5),\ (3,3)\}$.\\
		Znaleźć parę punktów w~najmniejszej odległości (euklidesowej).
	\end{example}

% ════════════════════════════════════════════════════════════════
\section{Mnożenie ciągu macierzy (Matrix-Chain-Order)}
% ════════════════════════════════════════════════════════════════

\begin{example}[Matrix-Chain-Order -- instancja 1]
	$n = 4$, \quad $p = (5,\ 4,\ 6,\ 2,\ 7)$, czyli\\
	$A_1\ (5\times 4),\ A_2\ (4\times 6),\ A_3\ (6\times 2),\ A_4\ (2\times 7)$.\\
	Wyznaczyć tablice $m$ i~$s$ oraz optymalne nawiasowanie.
\end{example}

\begin{example}[Matrix-Chain-Order -- instancja 2]
	$n = 5$, \quad $p = (10,\ 3,\ 12,\ 5,\ 50,\ 6)$, czyli\\
	$A_1\ (10\times 3),\ A_2\ (3\times 12),\ A_3\ (12\times 5),\ A_4\ (5\times 50),\ A_5\ (50\times 6)$.\\
	Wyznaczyć tablice $m$ i~$s$ oraz optymalne nawiasowanie.
\end{example}

	\begin{example}[Matrix-Chain-Order -- instancja 3]
		$n = 5$, \quad $p = (4,\ 10,\ 3,\ 12,\ 20,\ 7)$, czyli\\
		$A_1\ (4\times 10),\ A_2\ (10\times 3),\ A_3\ (3\times 12),\ A_4\ (12\times 20),\ A_5\ (20\times 7)$.\\
		Wyznaczyć tablice $m$ i~$s$ oraz optymalne nawiasowanie.
	\end{example}

% ════════════════════════════════════════════════════════════════
\section{Najdłuższy wspólny podciąg (LCS)}
% ════════════════════════════════════════════════════════════════

\begin{example}[LCS -- instancja 1]
	$X = (A,\ B,\ C,\ B,\ D,\ A,\ B)$, \quad $Y = (B,\ D,\ C,\ A,\ B,\ A)$.\\
	Wyznaczyć tablicę długości LCS oraz jeden najdłuższy wspólny podciąg.
\end{example}

\begin{example}[LCS -- instancja 2]
	$X = (G,\ A,\ T,\ T,\ A,\ C,\ A)$, \quad $Y = (G,\ C,\ A,\ T,\ A,\ A,\ T)$.\\
	Wyznaczyć tablicę długości LCS oraz jeden najdłuższy wspólny podciąg.
\end{example}

	\begin{example}[LCS -- instancja 3]
		$X = (A,\ B,\ A,\ C,\ B,\ D,\ A)$, \quad $Y = (B,\ A,\ B,\ C,\ D,\ A,\ B)$.\\
		Wyznaczyć tablicę długości LCS oraz jeden najdłuższy wspólny podciąg.
	\end{example}

% ════════════════════════════════════════════════════════════════
\section{Zachłanny wybór zajęć}
% ════════════════════════════════════════════════════════════════

\begin{example}[Wybór zajęć -- instancja 1]
	Znaleźć podzbiór o~największej liczności parami zgodnych zajęć:
	\begin{aztable}
		\centering
		\renewcommand{\arraystretch}{1.2}
		\begin{NiceTabular}{c|ccccccc}[margin,hvlines]
			$i$   & 1 & 2 & 3 & 4 & 5 & 6 & 7  \\
			\hline
			$s_i$ & 1 & 3 & 0 & 5 & 3 & 5 & 8  \\
			$t_i$ & 4 & 5 & 6 & 7 & 9 & 9 & 10
		\end{NiceTabular}
	\end{aztable}
\end{example}

\begin{example}[Wybór zajęć -- instancja 2]
	Znaleźć podzbiór o~największej liczności parami zgodnych zajęć:
	\begin{aztable}
		\centering
		\renewcommand{\arraystretch}{1.2}
		\begin{NiceTabular}{c|ccccccccc}[margin,hvlines]
			$i$   & 1 & 2 & 3 & 4 & 5 & 6 & 7 & 8  & 9  \\
			\hline
			$s_i$ & 0 & 2 & 1 & 6 & 4 & 8 & 7 & 9  & 11 \\
			$t_i$ & 3 & 5 & 4 & 8 & 7 & 9 & 11& 12 & 14
		\end{NiceTabular}
	\end{aztable}
\end{example}

	\begin{example}[Wybór zajęć -- instancja 3]
		Znaleźć podzbiór o~największej liczności parami zgodnych zajęć:
		\begin{aztable}
			\centering
			\renewcommand{\arraystretch}{1.2}
			\begin{NiceTabular}{c|cccccccc}[margin,hvlines]
				$i$   & 1 & 2 & 3 & 4 & 5 & 6 & 7 & 8  \\
				\hline
				$s_i$ & 1 & 2 & 0 & 3 & 4 & 6 & 5 & 7  \\
				$t_i$ & 3 & 5 & 7 & 8 & 9 & 10& 9 & 12
			\end{NiceTabular}
		\end{aztable}
	\end{example}

% ════════════════════════════════════════════════════════════════
\section{Algorytm Huffmana}
% ════════════════════════════════════════════════════════════════

\begin{example}[Huffman -- instancja 1]
	Znaki i~ich częstości:
	\begin{aztable}
		\centering
		\renewcommand{\arraystretch}{1.2}
		\begin{NiceTabular}{c|cccccc}[margin,hvlines]
			znak     & $a$ & $b$ & $c$ & $d$ & $e$ & $f$ \\
			\hline
			częstość & 45  & 13  & 12  & 16  & 9   & 5
		\end{NiceTabular}
	\end{aztable}
	Skonstruować drzewo Huffmana i~przypisać słowa kodowe.
\end{example}

\begin{example}[Huffman -- instancja 2]
	Zaczynamy ze zbiorem wierzchołków:
	$(g, 2),\ (o, 3),\ (k, 6),\ (m, 8),\ (l, 11),\ (u, 14),\ (e, 20)$.\\
	Skonstruować drzewo Huffmana i~przypisać słowa kodowe.
\end{example}

	\begin{example}[Huffman -- instancja 3]
		Znaki i~ich częstości:
		\begin{aztable}
			\centering
			\renewcommand{\arraystretch}{1.2}
			\begin{NiceTabular}{c|ccccccc}[margin,hvlines]
				znak     & $a$ & $b$ & $c$ & $d$ & $e$ & $f$ & $g$ \\
				\hline
				częstość & 30  & 25  & 20  & 15  & 10  & 7   & 3
			\end{NiceTabular}
		\end{aztable}
		Skonstruować drzewo Huffmana i~przypisać słowa kodowe.
	\end{example}

% ════════════════════════════════════════════════════════════════
\section{Szeregowanie zadań z deadlinami}
% ════════════════════════════════════════════════════════════════

\begin{example}[Szeregowanie zadań -- instancja 1]
	Zadania (kary $w_i$ za spóźnienie, deadliny $d_i$):
	\begin{aztable}
		\centering
		\renewcommand{\arraystretch}{1.2}
		\begin{NiceTabular}{c|ccccccc}[margin,hvlines]
			$i$   & 1   & 2  & 3  & 4  & 5  & 6  & 7  \\
			\hline
			$d_i$ & 4   & 2  & 4  & 3  & 1  & 4  & 6  \\
			$w_i$ & 100 & 90 & 70 & 60 & 50 & 30 & 15
		\end{NiceTabular}
	\end{aztable}
	Wyznaczyć harmonogram minimalizujący sumę kar zadań spóźnionych.
\end{example}

\begin{example}[Szeregowanie zadań -- instancja 2]
	Wyznaczyć harmonogram minimalizujący sumę kar dla następujących zadań:
	\begin{aztable}
		\centering
		\renewcommand{\arraystretch}{1.2}
		\begin{NiceTabular}{c|ccccccc}[margin,hvlines]
			$i$   & 1  & 2  & 3  & 4  & 5  & 6  & 7  \\
			\hline
			$d_i$ & 1  & 1  & 3  & 2  & 5  & 4  & 3  \\
			$w_i$ & 60 & 55 & 45 & 35 & 25 & 20 & 10
		\end{NiceTabular}
	\end{aztable}
\end{example}

	\begin{example}[Szeregowanie zadań -- instancja 3]
		Zadania (kary $w_i$ za spóźnienie, deadliny $d_i$); wyznaczyć harmonogram minimalizujący sumę kar:
		\begin{aztable}
			\centering
			\renewcommand{\arraystretch}{1.2}
			\begin{NiceTabular}{c|ccccccc}[margin,hvlines]
				$i$   & 1  & 2  & 3  & 4  & 5  & 6  & 7  \\
				\hline
				$d_i$ & 3  & 1  & 2  & 4  & 2  & 1  & 5  \\
				$w_i$ & 40 & 50 & 30 & 25 & 20 & 35 & 10
			\end{NiceTabular}
		\end{aztable}
	\end{example}

% ════════════════════════════════════════════════════════════════
\section{Greedy-Set-Cover}
% ════════════════════════════════════════════════════════════════

\begin{example}[Greedy Set Cover -- instancja 1]
	$X = \{1,2,3,4,5,6,7,8,9,10,11,12\}$,
	\begin{flalign*}
		S_1 &= \{1,2,3,4,5,6\},   & S_2 &= \{5,6,8,9\},      & S_3 &= \{1,4,7,10\},       & \\
		S_4 &= \{2,5,7,8,11\},    & S_5 &= \{3,6,9,12\},     & S_6 &= \{10,11\}.          &
	\end{flalign*}
	$\mathcal{F} = \{S_1,\dots,S_6\}$. Znaleźć pokrycie zwracane przez algorytm zachłanny.
\end{example}

\begin{example}[Greedy Set Cover -- instancja 2]
	$X = \{1,\dots,9\}$,
	\begin{flalign*}
		S_1 &= \{1,2,3,4,5\}, & S_2 &= \{6,7,8,9\}, & S_3 &= \{1,4,7\}, & \\
		S_4 &= \{2,5,8\},     & S_5 &= \{3,6,9\}.   &     &              &
	\end{flalign*}
	$\mathcal{F} = \{S_1,\dots,S_5\}$. Znaleźć pokrycie zwracane przez algorytm zachłanny.
\end{example}

	\begin{example}[Greedy Set Cover -- instancja 3]
		$X = \{1,2,3,4,5,6,7,8,9,10\}$,
		\begin{flalign*}
			S_1 &= \{1,2,3,4,5\}, & S_2 &= \{4,5,6,7\},    & S_3 &= \{1,6,8\},      & \\
			S_4 &= \{2,7,9\},     & S_5 &= \{3,8,9,10\},   & S_6 &= \{5,10\}.       &
		\end{flalign*}
		$\mathcal{F} = \{S_1,\dots,S_6\}$. Znaleźć pokrycie zwracane przez algorytm zachłanny.
	\end{example}

% ════════════════════════════════════════════════════════════════
\section{Approx-Vertex-Cover}
% ════════════════════════════════════════════════════════════════

\begin{example}[Approx Vertex Cover -- instancja 1]
	Wyznaczyć pokrycie wierzchołkowe zwracane przez Approx-Vertex-Cover dla grafu:
	\begin{azfigure}
		\centering
		\begin{tikzpicture}[every node/.style={circle, draw, minimum size=0.6cm, inner sep=1pt}]
			\node (a) at (0,2) {$a$};
			\node (b) at (2,2) {$b$};
			\node (c) at (4,2) {$c$};
			\node (d) at (0,0) {$d$};
			\node (e) at (2,0) {$e$};
			\node (f) at (4,0) {$f$};
			\draw (a)--(b); \draw (b)--(c); \draw (a)--(d);
			\draw (b)--(e); \draw (c)--(f); \draw (d)--(e);
			\draw (e)--(f); \draw (b)--(d);
		\end{tikzpicture}
	\end{azfigure}
\end{example}

\begin{example}[Approx Vertex Cover -- instancja 2]
	Wyznaczyć pokrycie wierzchołkowe zwracane przez Approx-Vertex-Cover dla grafu:
	\begin{azfigure}
		\centering
		\begin{tikzpicture}[every node/.style={circle, draw, minimum size=0.6cm, inner sep=1pt}]
			\node (a) at (0,2)   {$a$};
			\node (b) at (1.5,2) {$b$};
			\node (c) at (3,2)   {$c$};
			\node (d) at (4.5,2) {$d$};
			\node (e) at (0.75,0){$e$};
			\node (f) at (2.25,0){$f$};
			\node (g) at (3.75,0){$g$};
			\draw (a)--(b); \draw (b)--(c); \draw (c)--(d);
			\draw (a)--(e); \draw (b)--(e); \draw (b)--(f);
			\draw (c)--(f); \draw (c)--(g); \draw (d)--(g); \draw (e)--(f);
		\end{tikzpicture}
	\end{azfigure}
\end{example}

	\begin{example}[Approx Vertex Cover -- instancja 3]
		Wyznaczyć pokrycie wierzchołkowe zwracane przez Approx-Vertex-Cover dla grafu:
		\begin{azfigure}
			\centering
			\begin{tikzpicture}[every node/.style={circle, draw, minimum size=0.6cm, inner sep=1pt}]
				\node (a) at (0,2)   {$a$};
				\node (b) at (2,2)   {$b$};
				\node (c) at (4,2)   {$c$};
				\node (d) at (1,0.7) {$d$};
				\node (e) at (3,0.7) {$e$};
				\node (f) at (0,-0.6){$f$};
				\node (g) at (4,-0.6){$g$};
				\draw (a)--(b); \draw (b)--(c); \draw (a)--(d);
				\draw (b)--(d); \draw (b)--(e); \draw (c)--(e);
				\draw (d)--(f); \draw (e)--(g); \draw (d)--(e); \draw (f)--(g);
			\end{tikzpicture}
		\end{azfigure}
	\end{example}

% ════════════════════════════════════════════════════════════════
\section{Exact-Subset-Sum}
% ════════════════════════════════════════════════════════════════

\begin{example}[Exact Subset Sum -- instancja 1]
	$S = \{1,\ 4,\ 5,\ 11,\ 12\}$ (w~kolejności $x_1,\dots,x_5$), \quad $t = 21$.\\
	Wyznaczyć kolejne listy $L_i$ i~największą sumę podzbioru $\leq t$.
\end{example}

\begin{example}[Exact Subset Sum -- instancja 2]
	$S = \{2,\ 3,\ 7,\ 8,\ 10\}$ (w~kolejności $x_1,\dots,x_5$), \quad $t = 15$.\\
	Wyznaczyć kolejne listy $L_i$ i~największą sumę podzbioru $\leq t$.
\end{example}

	\begin{example}[Exact Subset Sum -- instancja 3]
		$S = \{3,\ 5,\ 6,\ 9,\ 11\}$ (w~kolejności $x_1,\dots,x_5$), \quad $t = 20$.\\
		Wyznaczyć kolejne listy $L_i$ i~największą sumę podzbioru $\leq t$.
	\end{example}

% ════════════════════════════════════════════════════════════════
\section{Approx-Subset-Sum}
% ════════════════════════════════════════════════════════════════

\begin{example}[Approx Subset Sum -- instancja 1]
	$S = \{106,\ 103,\ 205,\ 104\}$ (w~kolejności $x_1,\dots,x_4$), \quad $t = 312$, \quad $\varepsilon = 0{,}4$.\\
	Prześledzić działanie z~przycinaniem ($\delta = \tfrac{\varepsilon}{2n}$): wyznaczyć listy $L_i$ i~zwracaną sumę.
\end{example}

\begin{example}[Approx Subset Sum -- instancja 2]
	$S = \{10,\ 12,\ 17,\ 25\}$ (w~kolejności $x_1,\dots,x_4$), \quad $t = 50$, \quad $\varepsilon = 0{,}5$.\\
	Prześledzić działanie z~przycinaniem ($\delta = \tfrac{\varepsilon}{2n}$): wyznaczyć listy $L_i$ i~zwracaną sumę.
\end{example}

	\begin{example}[Approx Subset Sum -- instancja 3]
		$S = \{8,\ 15,\ 23,\ 42\}$ (w~kolejności $x_1,\dots,x_4$), \quad $t = 70$, \quad $\varepsilon = 0{,}3$.\\
		Prześledzić działanie z~przycinaniem ($\delta = \tfrac{\varepsilon}{2n}$): wyznaczyć listy $L_i$ i~zwracaną sumę.
	\end{example}

% ════════════════════════════════════════════════════════════════
\section{Maksymalne skojarzenie w grafie dwudzielnym (MM)}
% ════════════════════════════════════════════════════════════════

\begin{example}[Maximum Matching -- instancja 1]
	Wyznaczyć najliczniejsze skojarzenie (ścieżki powiększające) w~grafie:
	\begin{azfigure}
		\centering
		\begin{tikzpicture}[x=1.1cm, y=1.1cm]
			\foreach \x/\name in {0/a, 1/b, 2/c, 3/d} {
					\node[circle, fill=black, inner sep=1.5pt, label=above:$\name$] (u\name) at (\x,1) {};
				}
			\foreach \x/\name in {0/p, 1/q, 2/r, 3/s} {
					\node[circle, fill=black, inner sep=1.5pt, label=below:$\name$] (l\name) at (\x,0) {};
				}
			\node[right] at (3.3,1) {$V_1$};
			\node[right] at (3.3,0) {$V_2$};
			\draw (ua)--(lp); \draw (ua)--(lq);
			\draw (ub)--(lp); \draw (ub)--(lr);
			\draw (uc)--(lq); \draw (uc)--(lr); \draw (uc)--(ls);
			\draw (ud)--(ls);
		\end{tikzpicture}
	\end{azfigure}
\end{example}

\begin{example}[Maximum Matching -- instancja 2]
	Wyznaczyć najliczniejsze skojarzenie (ścieżki powiększające) w~grafie:
	\begin{azfigure}
		\centering
		\begin{tikzpicture}[x=1.0cm, y=1.1cm]
			\foreach \x/\name in {0/a, 1/b, 2/c, 3/d, 4/e} {
					\node[circle, fill=black, inner sep=1.5pt, label=above:$\name$] (u\name) at (\x,1) {};
				}
			\foreach \x/\name in {0/p, 1/q, 2/r, 3/s, 4/t} {
					\node[circle, fill=black, inner sep=1.5pt, label=below:$\name$] (l\name) at (\x,0) {};
				}
			\node[right] at (4.3,1) {$V_1$};
			\node[right] at (4.3,0) {$V_2$};
			\draw (ua)--(lp); \draw (ua)--(lq);
			\draw (ub)--(lp);
			\draw (uc)--(lq); \draw (uc)--(lr);
			\draw (ud)--(lr); \draw (ud)--(ls); \draw (ud)--(lt);
			\draw (ue)--(ls);
		\end{tikzpicture}
	\end{azfigure}
\end{example}

	\begin{example}[Maximum Matching -- instancja 3]
		Wyznaczyć najliczniejsze skojarzenie (ścieżki powiększające) w~grafie:
		\begin{azfigure}
			\centering
			\begin{tikzpicture}[x=1.1cm, y=1.1cm]
				\foreach \x/\name in {0/a, 1/b, 2/c, 3/d} {
						\node[circle, fill=black, inner sep=1.5pt, label=above:$\name$] (u\name) at (\x,1) {};
					}
				\foreach \x/\name in {0/p, 1/q, 2/r, 3/s} {
						\node[circle, fill=black, inner sep=1.5pt, label=below:$\name$] (l\name) at (\x,0) {};
					}
				\node[right] at (3.3,1) {$V_1$};
				\node[right] at (3.3,0) {$V_2$};
				\draw (ua)--(lp); \draw (ua)--(lr);
				\draw (ub)--(lp); \draw (ub)--(lq);
				\draw (uc)--(lq); \draw (uc)--(ls);
				\draw (ud)--(lr); \draw (ud)--(ls);
			\end{tikzpicture}
		\end{azfigure}
	\end{example}

% ════════════════════════════════════════════════════════════════
\appendix
\newpage
\section{Rozwiązania}
% ════════════════════════════════════════════════════════════════
% Każde rozwiązanie w osobnym pliku w rozwiazania/.
% ── Rozwiązanie: Para najbliższych punktów -- instancja 1 ──────────
\subsection{Para najbliższych punktów -- instancja 1}

% Kolory grup w drzewie rekursji (4 liście = 4 wywołania siłowe).
\colorlet{grpLL}{accent}%  {p8,p1}
\colorlet{grpLR}{accentC}% {p6,p4}
\colorlet{grpRL}{accentB}% {p7,p5}
\colorlet{grpRR}{violet}%  {p2,p3}
\colorlet{grpL}{accent}%   lewa połowa P_L
\colorlet{grpR}{accentB}%  prawa połowa P_R

% Pudełko grupujące krok rekursji; #1 = kolor (jak w drzewie).
\newtcolorbox{grpbox}[1][black]{%
	colback=#1!7, colframe=#1!55, boxrule=0.6pt,
	left=2mm, right=2mm, top=1mm, bottom=1mm, breakable}

Dane: $Q = \{(2,3),\ (12,30),\ (40,50),\ (5,1),\ (12,10),\ (3,4),\ (8,9),\ (1,1)\}$,
$n = 8$. Ponumerujmy punkty w~kolejności wejściowej:
\[
	p_1=(2,3),\ p_2=(12,30),\ p_3=(40,50),\ p_4=(5,1),\
	p_5=(12,10),\ p_6=(3,4),\ p_7=(8,9),\ p_8=(1,1).
\]

\paragraph{Sortowanie wstępne.}
Tworzymy tablice $X$ (rosnąco po~$x$) i~$Y$ (rosnąco po~$y$):
\begin{align*}
	X &= \big[\,p_8(1,1),\ p_1(2,3),\ p_6(3,4),\ p_4(5,1),\ p_7(8,9),\ p_5(12,10),\ p_2(12,30),\ p_3(40,50)\,\big], \\
	Y &= \big[\,p_8(1,1),\ p_4(5,1),\ p_1(2,3),\ p_6(3,4),\ p_7(8,9),\ p_5(12,10),\ p_2(12,30),\ p_3(40,50)\,\big].
\end{align*}

\paragraph{Dziel.}
$|P|=8>3$, dzielimy prostą $l\colon x=5$ na $|P_L|=|P_R|=4$, a~każdą połowę dalej, aż
do podzbiorów $\leq 3$ punktów (liście --- sprawdzane siłowo). Strukturę rekursji
i~kolory grup zbiera rysunek~\ref{fig:closest1:tree}.

\begin{azfigure}
	\centering
	\begin{tikzpicture}[
			font=\small,
			level distance=12mm,
			level 1/.style={sibling distance=42mm},
			level 2/.style={sibling distance=22mm},
			grp/.style={draw, rounded corners=2pt, inner sep=3pt, thick},
			edge from parent/.style={draw, gray, -}]
		\node[grp] {$P$ \scriptsize(8)}
			child {node[grp, grpL, text=grpL] {$P_L$ \scriptsize(4)}
				child {node[grp, grpLL, text=grpLL] {$\{p_8,p_1\}$}}
				child {node[grp, grpLR, text=grpLR] {$\{p_6,p_4\}$}}}
			child {node[grp, grpR, text=grpR] {$P_R$ \scriptsize(4)}
				child {node[grp, grpRL, text=grpRL] {$\{p_7,p_5\}$}}
				child {node[grp, grpRR, text=grpRR] {$\{p_2,p_3\}$}}};
	\end{tikzpicture}
	\caption{Drzewo rekursji. Cztery liście (po~2 punkty) liczone są siłowo;
		ich kolory powtarzają się w~tekście i~na rysunku~\ref{fig:closest1:result}.}
	\label{fig:closest1:tree}
\end{azfigure}

\[
	P_L=\{\textcolor{grpLL}{p_8,p_1},\ \textcolor{grpLR}{p_6,p_4}\},\qquad
	P_R=\{\textcolor{grpRL}{p_7,p_5},\ \textcolor{grpRR}{p_2,p_3}\}.
\]

\begin{grpbox}[grpL]
\paragraph{Zwyciężaj --- lewa połowa $P_L$.}
Liście (po~2 punkty --- siłowo):
\[
	d(\textcolor{grpLL}{p_8,p_1})=\sqrt{1^2+2^2}=\sqrt5\approx2{,}24,\qquad
	d(\textcolor{grpLR}{p_6,p_4})=\sqrt{2^2+3^2}=\sqrt{13}\approx3{,}61,
\]
więc $\delta=\sqrt5$. Faza \emph{połącz} w~$P_L$: pas $x\in[\,2-\sqrt5,\ 2+\sqrt5\,]$;
ponieważ $2<\sqrt5<3$, jego końce leżą w~$(-1,0)$ oraz $(4,5)$, więc pas
odrzuca $p_4(5,1)$, zostaje $Y'=[\,p_8(1,1),\,p_1(2,3),\,p_6(3,4)\,]$. Sprawdzamy:
\[
	d(p_8,p_1)=\sqrt5,\quad d(p_8,p_6)=\sqrt{13},\quad
	\boxed{d(p_1,p_6)=\sqrt{1^2+1^2}=\sqrt2\approx1{,}41.}
\]
Zatem $\delta_L=\sqrt2$ dla pary $\big(p_1(2,3),\,p_6(3,4)\big)$.
\end{grpbox}

\begin{grpbox}[grpR]
\paragraph{Zwyciężaj --- prawa połowa $P_R$.}
Liście:
\[
	d(\textcolor{grpRL}{p_7,p_5})=\sqrt{4^2+1^2}=\sqrt{17}\approx4{,}12,\qquad
	d(\textcolor{grpRR}{p_2,p_3})=\sqrt{28^2+20^2}=\sqrt{1184}\approx34{,}4,
\]
więc $\delta=\sqrt{17}$. Faza \emph{połącz}: ponieważ $4<\sqrt{17}<5$, pas
$x\in[\,12-\sqrt{17},\ 12+\sqrt{17}\,]$ ma końce w~$(7,8)$ oraz $(16,17)$, więc
$Y'=[\,p_7(8,9),\,p_5(12,10),\,p_2(12,30)\,]$ (odpada $p_3$). Żadna para z~$Y'$ nie
jest bliższa niż $\sqrt{17}$. Zatem $\delta_R=\sqrt{17}$ dla pary $\big(p_7(8,9),\,p_5(12,10)\big)$.
\end{grpbox}

\paragraph{Połącz --- poziom górny.}
$\delta=\min(\delta_L,\delta_R)=\min(\sqrt2,\sqrt{17})=\sqrt2$.
Pas wokół $l\colon x=5$ ma szerokość $2\delta=2\sqrt2$; ponieważ $1<\sqrt2<2$,
pas $x\in[\,5-\sqrt2,\ 5+\sqrt2\,]$ ma końce w~$(3,4)$ oraz $(6,7)$. W~pasie leży tylko
$p_4(5,1)$, więc $Y'=[\,p_4\,]$ --- brak pary „międzystronowej”, $\delta$ bez zmian.

\paragraph{Wynik.}
Para najbliższych punktów to $(2,3)$ i~$(3,4)$ w~odległości $\delta=\sqrt2\approx1{,}41$.

\begin{azfigure}
	\centering
	\begin{tikzpicture}[>=stealth, font=\small, scale=0.18]
		\draw[gray!50, very thin] (0,0) grid[step=5] (40,50);
		\draw[->] (0,0)--(42,0) node[right]{$x$};
		\draw[->] (0,0)--(0,52) node[above]{$y$};
		\draw[thick, accent] (5,-2)--(5,52) node[above]{$l$};
		% punkty pokolorowane wg liścia drzewa rekursji
		\foreach \x/\y/\c in {1/1/grpLL,2/3/grpLL,3/4/grpLR,5/1/grpLR,%
				8/9/grpRL,12/10/grpRL,12/30/grpRR,40/50/grpRR}
			\fill[\c] (\x,\y) circle (7pt);
		% znaleziona para
		\draw[red, thick] (2,3)--(3,4);
		\draw[red, thick] (2,3) circle (12pt);
		\draw[red, thick] (3,4) circle (12pt);
		\node[red, below left=2pt, scale=1.5] at (2,3) {$(2,3)$};
		\node[red, right=4pt, scale=1.5] at (3,4) {$(3,4)$};
	\end{tikzpicture}
	\caption{Instancja~1: prosta podziału $l\colon x=5$, punkty pokolorowane wg
		liścia drzewa rekursji oraz znaleziona para $(2,3)$, $(3,4)$ (czerwona).}
	\label{fig:closest1:result}
\end{azfigure}

% ── Rozwiązanie: Matrix-Chain-Order -- instancja 1 ────────────────
\subsection{Matrix-Chain-Order -- instancja 1}

Dane: $n=4$, $p=(5,\ 4,\ 6,\ 2,\ 7)$, czyli
\[
	A_1\ (5\times 4),\quad A_2\ (4\times 6),\quad A_3\ (6\times 2),\quad A_4\ (2\times 7).
\]
Wypełniamy tablice $m$ i~$s$ \emph{przekątnymi} po~rosnącej długości
podciągu $L$; w~każdym kroku dopisane komórki wyróżniono na
\textcolor{accentB}{pomarańczowo}, a~$\times$ to pola nieużywane
($i>j$ dla~$m$, $i\geq j$ dla~$s$). Skrót: $A_i=p_{i-1}\times p_i$.

\medskip
\noindent Stan początkowy ($L=1$, pojedyncze macierze, koszt $0$):\\[0.7em]
\begin{minipage}{0.45\linewidth}
	$m$
	$\begin{bNiceArray}{c|cccc}[margin,hvlines]
			i \backslash j & 1 & 2 & 3 & 4 \\
			1 & 0 &  &  &  \\
			2 & \times & 0 &  &  \\
			3 & \times & \times & 0 &  \\
			4 & \times & \times & \times & 0
		\end{bNiceArray}$
\end{minipage}\hfill
\begin{minipage}{0.45\linewidth}
	$s$
	$\begin{bNiceArray}{c|cccc}[margin,hvlines]
			i \backslash j & 1 & 2 & 3 & 4 \\
			1 & \times &  &  &  \\
			2 & \times & \times &  &  \\
			3 & \times & \times & \times &  \\
			4 & \times & \times & \times & \times
		\end{bNiceArray}$
\end{minipage}

\medskip
\noindent{\color{placeholdercolor}\rule{\linewidth}{0.4pt}}\par\smallskip
\noindent{\color{dygresjagray}\footnotesize$\begin{array}{c|*{5}{c}} i & 0 & 1 & 2 & 3 & 4 \\ \hline p_i & 5 & 4 & 6 & 2 & 7 \end{array}$\qquad $A_i = p_{i-1}\times p_i$}
\begin{flalign*}
	L=2: \quad & m(1, 2) = 5 \cdot 4 \cdot 6 = 120 &  & \\
	           & m(2, 3) = 4 \cdot 6 \cdot 2 = 48  &  & \\
	           & m(3, 4) = 6 \cdot 2 \cdot 7 = 84  &  &
\end{flalign*}

\medskip
\noindent po $L=2$ (pary sąsiednich macierzy):\\[0.7em]
\begin{minipage}{0.45\linewidth}
	$m$
	$\begin{bNiceArray}{c|cccc}[margin,hvlines]
			i \backslash j & 1 & 2 & 3 & 4 \\
			1 & 0 & \color{accentB}120 &  &  \\
			2 & \times & 0 & \color{accentB}48 &  \\
			3 & \times & \times & 0 & \color{accentB}84 \\
			4 & \times & \times & \times & 0
		\end{bNiceArray}$
\end{minipage}\hfill
\begin{minipage}{0.45\linewidth}
	$s$
	$\begin{bNiceArray}{c|cccc}[margin,hvlines]
			i \backslash j & 1 & 2 & 3 & 4 \\
			1 & \times & \color{accentB}1 &  &  \\
			2 & \times & \times & \color{accentB}2 &  \\
			3 & \times & \times & \times & \color{accentB}3 \\
			4 & \times & \times & \times & \times
		\end{bNiceArray}$
\end{minipage}

\medskip
\noindent{\color{placeholdercolor}\rule{\linewidth}{0.4pt}}\par\smallskip
\noindent{\color{dygresjagray}\footnotesize$\begin{array}{c|*{5}{c}} i & 0 & 1 & 2 & 3 & 4 \\ \hline p_i & 5 & 4 & 6 & 2 & 7 \end{array}$\qquad $A_i = p_{i-1}\times p_i$}
\begin{flalign*}
	L=3: \quad & m(1, 3) = \min \begin{cases}
		                            m(1, 1) + m(2, 3) + 5 \cdot 4 \cdot 2 \\
		                            m(1, 2) + m(3, 3) + 5 \cdot 6 \cdot 2
	                            \end{cases} = \min \begin{cases}
		                                               0 + 48 + 40 \\
		                                               120 + 0 + 60
	                                               \end{cases} = \min \begin{cases} 88 \\ 180 \end{cases} = 88 \\
	           & m(2, 4) = \min \begin{cases}
		                            m(2, 2) + m(3, 4) + 4 \cdot 6 \cdot 7 \\
		                            m(2, 3) + m(4, 4) + 4 \cdot 2 \cdot 7
	                            \end{cases} = \min \begin{cases}
		                                               0 + 84 + 168 \\
		                                               48 + 0 + 56
	                                               \end{cases} = \min \begin{cases} 252 \\ 104 \end{cases} = 104
\end{flalign*}

\medskip
\noindent po $L=3$:\\[0.7em]
\begin{minipage}{0.45\linewidth}
	$m$
	$\begin{bNiceArray}{c|cccc}[margin,hvlines]
			i \backslash j & 1 & 2 & 3 & 4 \\
			1 & 0 & 120 & \color{accentB}88 &  \\
			2 & \times & 0 & 48 & \color{accentB}104 \\
			3 & \times & \times & 0 & 84 \\
			4 & \times & \times & \times & 0
		\end{bNiceArray}$
\end{minipage}\hfill
\begin{minipage}{0.45\linewidth}
	$s$
	$\begin{bNiceArray}{c|cccc}[margin,hvlines]
			i \backslash j & 1 & 2 & 3 & 4 \\
			1 & \times & 1 & \color{accentB}1 &  \\
			2 & \times & \times & 2 & \color{accentB}3 \\
			3 & \times & \times & \times & 3 \\
			4 & \times & \times & \times & \times
		\end{bNiceArray}$
\end{minipage}

\medskip
\noindent{\color{placeholdercolor}\rule{\linewidth}{0.4pt}}\par\smallskip
\noindent{\color{dygresjagray}\footnotesize$\begin{array}{c|*{5}{c}} i & 0 & 1 & 2 & 3 & 4 \\ \hline p_i & 5 & 4 & 6 & 2 & 7 \end{array}$\qquad $A_i = p_{i-1}\times p_i$}
\begin{flalign*}
	L=4: \quad & m(1, 4) = \min \begin{cases}
		                            m(1, 1) + m(2, 4) + 5 \cdot 4 \cdot 7 \\
		                            m(1, 2) + m(3, 4) + 5 \cdot 6 \cdot 7 \\
		                            m(1, 3) + m(4, 4) + 5 \cdot 2 \cdot 7
	                            \end{cases} = \min \begin{cases}
		                                               0 + 104 + 140 \\
		                                               120 + 84 + 210 \\
		                                               88 + 0 + 70
	                                               \end{cases} = \min \begin{cases} 244 \\ 414 \\ 158 \end{cases} = 158
\end{flalign*}

\medskip
\noindent tablice końcowe:\\[0.7em]
\begin{minipage}{0.45\linewidth}
	$m$
	$\begin{bNiceArray}{c|cccc}[margin,hvlines]
			i \backslash j & 1 & 2 & 3 & 4 \\
			1 & 0 & 120 & 88 & \color{accentB}158 \\
			2 & \times & 0 & 48 & 104 \\
			3 & \times & \times & 0 & 84 \\
			4 & \times & \times & \times & 0
		\end{bNiceArray}$
\end{minipage}\hfill
\begin{minipage}{0.45\linewidth}
	$s$
	$\begin{bNiceArray}{c|cccc}[margin,hvlines]
			i \backslash j & 1 & 2 & 3 & 4 \\
			1 & \times & 1 & 1 & \color{accentB}3 \\
			2 & \times & \times & 2 & 3 \\
			3 & \times & \times & \times & 3 \\
			4 & \times & \times & \times & \times
		\end{bNiceArray}$
\end{minipage}

\paragraph{Optymalne nawiasowanie.}
Odtwarzamy je z~tablicy $s$, rekurencyjnie dzieląc $A_i\dots A_j$ w~miejscu
$s(i,j)$:
\[
	s(1,4)=3 \;\Rightarrow\; (A_1 A_2 A_3)(A_4),\qquad
	s(1,3)=1 \;\Rightarrow\; (A_1)(A_2 A_3).
\]
Stąd optymalne nawiasowanie
\[
	\boxed{\big(A_1\,(A_2\,A_3)\big)\,A_4}
\]
o~minimalnym koszcie $m(1,4)=158$ mnożeń skalarnych.

% ── Rozwiązanie: LCS -- instancja 1 ───────────────────────────────
\subsection{Najdłuższy wspólny podciąg (LCS) -- instancja 1}

Dane: $X = (A,\ B,\ C,\ B,\ D,\ A,\ B)$, $Y = (B,\ D,\ C,\ A,\ B,\ A)$, czyli
$m = 7$, $n = 6$. Tablicę kosztów $c$ wypełniamy wierszami, korzystając z~reguły
\[
	c(i, j) = \begin{cases}
		          0                               & \text{jeśli } i = 0 \text{ lub } j = 0,         \\
		          c(i - 1, j - 1) + 1             & \text{jeśli } i, j > 0 \text{ i } x_i = y_j,    \\
		          \max\{c(i - 1, j),\ c(i, j - 1)\} & \text{jeśli } i, j > 0 \text{ i } x_i \neq y_j,
	          \end{cases}
\]
a~równolegle zapisujemy w~tablicy $b$ kierunek, z~którego wzięliśmy wartość
($\nwarrow$ -- dopasowanie, $\uparrow$ -- góra, $\leftarrow$ -- lewo; remisy na
korzyść góry). Komórki wchodzące do wybranego LCS wyróżniono na
\textcolor{accentB}{pomarańczowo}.

\begin{aztable}
	\centering
	\renewcommand{\arraystretch}{1.2}
	\begin{minipage}{0.48\linewidth}
		\centering
		Tablica kosztów $c$: \\[0.5em]
		\begin{NiceTabular}{cc|c|cccccc}[margin,hvlines]
			    & $j$   & 0     & 1   & 2   & 3   & 4   & 5   & 6   \\
			$i$ &       & $y_j$ & $B$ & $D$ & $C$ & $A$ & $B$ & $A$ \\
			\hline
			0   & $x_i$ & 0     & 0   & 0   & 0   & 0   & 0   & 0   \\
			1   & $A$   & 0     & 0   & 0   & 0   & 1   & 1   & 1   \\
			2   & $B$   & 0     & \color{accentB}1 & 1 & 1 & 1 & 2 & 2 \\
			3   & $C$   & 0     & 1   & 1   & \color{accentB}2 & 2 & 2 & 2 \\
			4   & $B$   & 0     & 1   & 1   & 2   & 2   & \color{accentB}3 & 3 \\
			5   & $D$   & 0     & 1   & 2   & 2   & 2   & 3   & 3   \\
			6   & $A$   & 0     & 1   & 2   & 2   & 3   & 3   & \color{accentB}4 \\
			7   & $B$   & 0     & 1   & 2   & 2   & 3   & 4   & 4   \\
		\end{NiceTabular}
	\end{minipage}\hfill
	\begin{minipage}{0.48\linewidth}
		\centering
		Tablica strzałek $b$: \\[0.5em]
		\begin{NiceTabular}{cc|c|cccccc}[margin,hvlines]
			    & $j$   & 0     & 1          & 2            & 3            & 4            & 5            & 6            \\
			$i$ &       & $y_j$ & $B$        & $D$          & $C$          & $A$          & $B$          & $A$          \\
			\hline
			0   & $x_i$ &       &            &              &              &              &              &              \\
			1   & $A$   &       & $\uparrow$ & $\uparrow$   & $\uparrow$   & $\nwarrow$   & $\leftarrow$ & $\nwarrow$   \\
			2   & $B$   &       & \color{accentB}$\nwarrow$ & $\leftarrow$ & $\leftarrow$ & $\uparrow$   & $\nwarrow$   & $\leftarrow$ \\
			3   & $C$   &       & $\uparrow$ & $\uparrow$   & \color{accentB}$\nwarrow$ & $\leftarrow$ & $\uparrow$   & $\uparrow$   \\
			4   & $B$   &       & $\nwarrow$ & $\uparrow$   & $\uparrow$   & $\uparrow$   & \color{accentB}$\nwarrow$ & $\leftarrow$ \\
			5   & $D$   &       & $\uparrow$ & $\nwarrow$   & $\uparrow$   & $\uparrow$   & $\uparrow$   & $\uparrow$   \\
			6   & $A$   &       & $\uparrow$ & $\uparrow$   & $\uparrow$   & $\nwarrow$   & $\uparrow$   & \color{accentB}$\nwarrow$ \\
			7   & $B$   &       & $\nwarrow$ & $\uparrow$   & $\uparrow$   & $\uparrow$   & $\nwarrow$   & $\uparrow$   \\
		\end{NiceTabular}
	\end{minipage}
\end{aztable}

\begin{dygresja}
	Po jednej komórce na każdą gałąź reguły (pozostałe $42$ liczymy analogicznie),
	$X = (x_1, \dots, x_7)$, $Y = (y_1, \dots, y_6)$:
	\begin{itemize}
		\item \textbf{Dopasowanie} ($\nwarrow$): $c(2, 1)$, $x_2 = B = y_1$, więc
		      $c(2, 1) = c(1, 0) + 1 = 1$.
		\item \textbf{Góra} ($\uparrow$): $c(3, 1)$, $x_3 = C \neq B = y_1$; sąsiedzi
		      $c(2, 1) = 1 \geq c(3, 0) = 0$, więc $c(3, 1) = 1$.
		\item \textbf{Lewo} ($\leftarrow$): $c(2, 2)$, $x_2 = B \neq D = y_2$; sąsiedzi
		      $c(1, 2) = 0 < c(2, 1) = 1$, więc $c(2, 2) = 1$.
		\item \textbf{Remis}: $c(1, 1)$, $A \neq B$, $c(0, 1) = c(1, 0) = 0$ -- warunek
		      $\geq$ spełniony, więc $\uparrow$.
	\end{itemize}
\end{dygresja}

\paragraph{Odtworzenie podciągu.}
Startujemy z~$(m, n) = (7, 6)$ i~idziemy za strzałkami $b$ aż do brzegu; przy
$\nwarrow$ schodzimy po~przekątnej, dopisując $x_i$ do~LCS:
\[
	\begin{aligned}
		(7, 6) & \xrightarrow{\uparrow} (6, 6) \xrightarrow{\nwarrow} (5, 5) \;[\,x_6 = A\,]
		\xrightarrow{\uparrow} (4, 5) \xrightarrow{\nwarrow} (3, 4) \;[\,x_4 = B\,] \\
		       & \xrightarrow{\leftarrow} (3, 3) \xrightarrow{\nwarrow} (2, 2) \;[\,x_3 = C\,]
		\xrightarrow{\leftarrow} (2, 1) \xrightarrow{\nwarrow} (1, 0) \;[\,x_2 = B\,].
	\end{aligned}
\]
Znaki dopisywane są \emph{po} powrocie z~rekurencji, więc w~kolejności od
najgłębszego: $B, C, B, A$. Stąd najdłuższy wspólny podciąg
\[
	\boxed{\text{LCS}(X, Y) = (B,\ C,\ B,\ A)}
\]
o~długości $c(7, 6) = 4$. (Nie jest jedyny -- np. $(B, D, A, B)$ też ma długość~$4$.)

% ── Rozwiązanie: Zachłanny wybór zajęć -- instancja 1 ─────────────
\subsection{Zachłanny wybór zajęć -- instancja 1}

Dane ($n = 7$ zajęć), już posortowane po~rosnącym czasie zakończenia $t_i$.
Algorytm \textsc{GreedyActivitySelector} wybiera $z_1$, a~następnie idzie po
zajęciach rosnąco i~dodaje $z_i$ do~rozwiązania $A$ wtedy i~tylko wtedy, gdy jest
zgodne z~ostatnio wybranym $z_j$, tj. $s_i \geq t_j$. Wybrane zajęcia wyróżniono
na~\textcolor{accentB}{pomarańczowo}.

\begin{aztable}
	\centering
	\renewcommand{\arraystretch}{1.2}
	\begin{NiceTabular}{c|ccccccc}[margin,hvlines]
		$i$   & \color{accentB}1 & 2 & 3 & \color{accentB}4 & 5 & 6 & \color{accentB}7  \\
		\hline
		$s_i$ & \color{accentB}1 & 3 & 0 & \color{accentB}5 & 3 & 5 & \color{accentB}8  \\
		$t_i$ & \color{accentB}4 & 5 & 6 & \color{accentB}7 & 9 & 9 & \color{accentB}10
	\end{NiceTabular}
\end{aztable}

\begin{dygresja}
	Przebieg pętli ($j$ -- indeks ostatnio wybranego zajęcia, początkowo
	$A = \{z_1\}$, $j = 1$, $t_j = 4$):
	\begin{itemize}
		\item $i = 2$: $s_2 = 3 \geq t_j = 4$? \emph{nie} -- pomijamy.
		\item $i = 3$: $s_3 = 0 \geq 4$? \emph{nie} -- pomijamy.
		\item $i = 4$: $s_4 = 5 \geq 4$? \emph{tak} -- $A \cup \{z_4\}$, $j \Assign 4$, $t_j = 7$.
		\item $i = 5$: $s_5 = 3 \geq 7$? \emph{nie} -- pomijamy.
		\item $i = 6$: $s_6 = 5 \geq 7$? \emph{nie} -- pomijamy.
		\item $i = 7$: $s_7 = 8 \geq 7$? \emph{tak} -- $A \cup \{z_7\}$, $j \Assign 7$.
	\end{itemize}
\end{dygresja}

\paragraph{Wynik.}
Wybrane zajęcia są parami zgodne, bo ich przedziały
$\langle 1, 4)$, $\langle 5, 7)$, $\langle 8, 10)$ są rozłączne:
\[
	\boxed{A = \{z_1,\ z_4,\ z_7\}}
\]
o~liczności $|A| = 3$. 

% ── Rozwiązanie: Huffman -- instancja 1 ──────────────────────────
\subsection{Algorytm Huffmana -- instancja 1}

Dane: $C = \{a, b, c, d, e, f\}$ z~częstościami
$f(a) = 45,\ f(b) = 13,\ f(c) = 12,\ f(d) = 16,\ f(e) = 9,\ f(f) = 5$.
Algorytm Huffmana buduje drzewo \emph{od dołu}: w~każdej z~$n - 1 = 5$ iteracji
wyjmuje z~kolejki dwa korzenie o~\emph{najmniejszej} częstości, łączy je nowym
wierzchołkiem $z$ o~częstości $f(z) = f(x) + f(y)$ (mniejszy jako lewe dziecko)
i~wkłada $z$ z~powrotem. Krawędź do lewego dziecka etykietujemy $0$, do~prawego $1$.

% Styl drzew jak w wykładzie (lect07): okrągłe, jednakowe węzły z etykietą "a,45".
\tikzset{
	htree/.style={
			level distance=0.95cm,
			level 1/.style={sibling distance=4cm},
			level 2/.style={sibling distance=2.5cm},
			level 3/.style={sibling distance=1.5cm},
			level 4/.style={sibling distance=1.0cm},
			every node/.style={circle, draw, minimum size=0.7cm, inner sep=0pt, font=\footnotesize, align=center},
			elabel left/.style={draw=none, font=\scriptsize, inner sep=1pt, midway, sloped, above},
			elabel right/.style={draw=none, font=\scriptsize, inner sep=1pt, midway, sloped, above},
		}
}
\newcommand{\huffstep}[2]{%
	\noindent
	\begin{minipage}[t]{0.34\textwidth}\vspace{0pt}\raggedright #1\end{minipage}%
	\hfill
	\begin{minipage}[t]{0.62\textwidth}\vspace{0pt}\centering#2\end{minipage}%
	\par\vspace{2em}%
}

\huffstep{%
	\textbf{1.} Łączymy $f$ i~$e$. \\
	Zostają: $(c, 12), (b, 13), (14), (d, 16), (a, 45)$.
}{%
	\begin{tikzpicture}[htree]
		\node {14}
		child { node {f,5} }
		child { node {e,9} };
	\end{tikzpicture}
}

\huffstep{%
	\textbf{2.} Łączymy $c$ i~$b$. \\
	Zostają: $(14), (d, 16), (25), (a, 45)$.
}{%
	\begin{tikzpicture}[htree]
		\node {25}
		child { node {c,12} }
		child { node {b,13} };
	\end{tikzpicture}
}

\huffstep{%
	\textbf{3.} Łączymy $14$ i~$d$. \\
	Zostają: $(25), (30), (a, 45)$.
}{%
	\begin{tikzpicture}[htree]
		\node {30}
		child { node {14}
				child { node {f,5} }
				child { node {e,9} }
			}
		child { node {d,16} };
	\end{tikzpicture}
}

\huffstep{%
	\textbf{4.} Łączymy $25$ i~$30$. \\
	Zostają: $(a, 45), (55)$.
}{%
	\begin{tikzpicture}[htree]
		\node {55}
		child { node {25}
				child { node {c,12} }
				child { node {b,13} }
			}
		child { node {30}
				child { node {14}
						child { node {f,5} }
						child { node {e,9} }
					}
				child { node {d,16} }
			};
	\end{tikzpicture}
}

\huffstep{%
	\textbf{5.} Łączymy $a$ i~$55$. \\
	(ostateczne drzewo kodów)
}{%
	\begin{tikzpicture}[htree]
		\node {100}
		child { node {a,45} edge from parent node[elabel left] {0} }
		child { node {55}
				child { node {25}
						child { node {c,12} edge from parent node[elabel left] {0} }
						child { node {b,13} edge from parent node[elabel right] {1} }
						edge from parent node[elabel left] {0}
					}
				child { node {30}
						child { node {14}
								child { node {f,5} edge from parent node[elabel left] {0} }
								child { node {e,9} edge from parent node[elabel right] {1} }
								edge from parent node[elabel left] {0}
							}
						child { node {d,16} edge from parent node[elabel right] {1} }
						edge from parent node[elabel right] {1}
					}
				edge from parent node[elabel right] {1}
			};
	\end{tikzpicture}
}

Odczytany kod z~drzewa w~kroku 5:
\begin{flalign*}
	a & \implies 0    &  & \\
	c & \implies 100  &  & \\
	b & \implies 101  &  & \\
	d & \implies 111  &  & \\
	f & \implies 1100 &  & \\
	e & \implies 1101 &  &
\end{flalign*}

\paragraph{Liczba bitów.}
Długość zakodowanego pliku to $B(T) = \sum_{c \in C} f(c)\, d_T(c)$:
\[
	\boxed{B(T) = 45\cdot 1 + 12\cdot 3 + 13\cdot 3 + 16\cdot 3 + 5\cdot 4 + 9\cdot 4 = 224 \text{ bitów}}
\]
(wobec $58 \cdot 3 = 300$ bitów przy kodzie stałej długości po $3$ bity na znak).

% ── Rozwiązanie: Szeregowanie zadań z deadlinami -- instancja 1 ────
\subsection{Szeregowanie zadań z deadlinami -- instancja 1}

Dane: $n = 7$ zadań o~jednostkowym czasie wykonania, deadliny $d_i$ i~kary
$w_i$. Rodzina zbiorów \emph{wolnych od kar} tworzy matroid, więc optymalne
rozwiązanie daje algorytm \textsc{Greedy}: bierzemy zadania w~kolejności
\emph{nierosnącej} kary i~dodajemy $z_i$ do~$A$ wtedy i~tylko wtedy, gdy $A$
pozostaje wolny od kar, tzn.\ profil $N_t(A) = |\{z_j \in A : d_j \leq t\}|$
nie przekracza $(1, 2, \ldots, n)$. Zadania wybrane (terminowe) wyróżniono na~%
\textcolor{accentB}{pomarańczowo}.

\begin{aztable}
	\centering
	\renewcommand{\arraystretch}{1.2}
	\begin{NiceTabular}{c|ccccccc}[margin,hvlines]
		$i$   & \color{accentB}1 & \color{accentB}2 & \color{accentB}3 & \color{accentB}4 & 5 & 6 & \color{accentB}7  \\
		\hline
		$d_i$ & \color{accentB}4 & \color{accentB}2 & \color{accentB}4 & \color{accentB}3 & 1 & 4 & \color{accentB}6  \\
		$w_i$ & \color{accentB}100 & \color{accentB}90 & \color{accentB}70 & \color{accentB}60 & 50 & 30 & \color{accentB}15
	\end{NiceTabular}
\end{aztable}

\begin{dygresja}
	Kary są już posortowane nierosnąco ($w_1 \geq \dots \geq w_7$), więc zadania
	rozważamy w~kolejności $z_1, \dots, z_7$. Poniżej $A$ wypisane rosnąco po
	deadlinach, a~po prawej profil $N(A) = (N_1, \ldots, N_7)$, który musi
	pozostać $\leq (1, 2, 3, 4, 5, 6, 7)$:
	\begin{enumerate}[nosep]
		\item $A = (z_1)$ \hfill $N(A) = (0, 0, 0, 1, 1, 1, 1)$
		\item $A = (z_2, z_1)$ \hfill $N(A) = (0, 1, 1, 2, 2, 2, 2)$
		\item $A = (z_2, z_1, z_3)$ \hfill $N(A) = (0, 1, 1, 3, 3, 3, 3)$
		\item $A = (z_2, z_4, z_1, z_3)$ \hfill $N(A) = (0, 1, 2, 4, 4, 4, 4)$
		\item $A = (z_2, z_4, z_1, z_3)$ \quad ($z_5$ nie dodajemy, bo $N_4(A \cup \{z_5\}) = 5 > 4$)
		\item $A = (z_2, z_4, z_1, z_3)$ \quad ($z_6$ nie dodajemy, bo $N_4(A \cup \{z_6\}) = 5 > 4$)
		\item $A = (z_2, z_4, z_1, z_3, z_7)$ \hfill $N(A) = (0, 1, 2, 4, 4, 5, 5)$
	\end{enumerate}
\end{dygresja}

\paragraph{Harmonogram.}
Zadania terminowe ustawiamy rosnąco po~deadlinach, a~spóźnione ($z_5, z_6$) na
końcu:
\[
	\boxed{(z_2,\ z_4,\ z_1,\ z_3,\ z_7,\ z_5,\ z_6)}
\]
Suma kar zadań spóźnionych: $w_5 + w_6 = 50 + 30 = 80$.

% ── Rozwiązanie: Greedy Set Cover -- instancja 1 ─────────────────
\subsection{Greedy Set Cover -- instancja 1}

Dane: $X = \{1, \ldots, 12\}$ oraz rodzina $\mathcal{F} = \{S_1, \ldots, S_6\}$.
Algorytm \textsc{GSC} startuje z~$U = X$ i~w~każdej iteracji dokłada do~pokrycia
$\mathcal{C}$ ten zbiór $S \in \mathcal{F}$, który pokrywa najwięcej elementów
jeszcze niepokrytych, tj.\ maksymalizuje $|S \cap U|$ (remisy rozstrzygamy na
korzyść mniejszego indeksu), po~czym odejmuje $U \Assign U \setminus S$. Zbiory
wybrane do~pokrycia wyróżniono na~\textcolor{accentB}{pomarańczowo}.

\begin{aztable}
	\centering
	\renewcommand{\arraystretch}{1.2}
	\begin{NiceTabular}{c|l|c}[margin,hvlines]
		$S$                  & elementy                & $|S|$ \\
		\hline
		\color{accentB}$S_1$ & \color{accentB}$\{1,2,3,4,5,6\}$ & \color{accentB}6 \\
		$S_2$                & $\{5,6,8,9\}$           & 4 \\
		\color{accentB}$S_3$ & \color{accentB}$\{1,4,7,10\}$    & \color{accentB}4 \\
		\color{accentB}$S_4$ & \color{accentB}$\{2,5,7,8,11\}$  & \color{accentB}5 \\
		\color{accentB}$S_5$ & \color{accentB}$\{3,6,9,12\}$    & \color{accentB}4 \\
		$S_6$                & $\{10,11\}$             & 2
	\end{NiceTabular}
\end{aztable}

\begin{dygresja}
	W~każdym wierszu po~lewej bieżący zbiór niepokrytych $U$, a~po~prawej
	rozważane przecięcia $|S \cap U|$ (pogrubiony zwycięzca):
	\begin{enumerate}[nosep]
		\item $U = \{1,\ldots,12\}$:\quad
		      $|S_1 \cap U| = \mathbf{6}$, $|S_4 \cap U| = 5$, pozostałe $\leq 4$
		      \hfill wybór $S_1$
		\item $U = \{7,8,9,10,11,12\}$:\quad
		      $|S_4 \cap U| = \mathbf{3}$, $|S_2 \cap U| = |S_3 \cap U| = |S_5 \cap U|
		      = |S_6 \cap U| = 2$
		      \hfill wybór $S_4$
		\item $U = \{9,10,12\}$:\quad
		      $|S_5 \cap U| = \mathbf{2}$, $|S_2 \cap U| = |S_3 \cap U|
		      = |S_6 \cap U| = 1$
		      \hfill wybór $S_5$
		\item $U = \{10\}$:\quad
		      $|S_3 \cap U| = |S_6 \cap U| = 1$ -- remis, mniejszy indeks
		      \hfill wybór $S_3$
	\end{enumerate}
	Po~czwartej iteracji $U = \emptyset$, pętla się kończy.
\end{dygresja}

\paragraph{Wynik.}
Algorytm zachłanny zwraca pokrycie złożone z~czterech zbiorów:
\[
	\boxed{\mathcal{C} = \{S_1,\ S_4,\ S_5,\ S_3\}}
\]
Nie jest ono optymalne: $\mathcal{C}^* = \{S_3, S_4, S_5\}$ pokrywa całe $X$
trzema zbiorami. Pierwszy, najobfitszy wybór $S_1$ okazał się pułapką --- po~nim
do~domknięcia pokrycia trzeba było aż trzech kolejnych zbiorów.

% ── Rozwiązanie: Approx Vertex Cover -- instancja 1 ──────────────
\subsection{Approx Vertex Cover -- instancja 1}

Algorytm \textsc{AVC} startuje z~$F = E(G)$ i~dopóki zostają niepokryte
krawędzie, wybiera dowolną krawędź $uv \in F$, dorzuca do~pokrycia $C$
\emph{oba} jej końce, a~następnie usuwa z~$F$ wszystkie krawędzie incydentne
z~$u$ lub $v$. Wybór krawędzi jest dowolny --- przyjmijmy kolejność
leksykograficzną najmniejszej pozostałej krawędzi. Krawędzie wybrane
(skojarzenie $B$) zaznaczono na~\textcolor{accentB}{pomarańczowo}, krawędzie
usuwane z~$F$ wraz z~nimi --- szaro, a~wierzchołki dokładane do~$C$ --- na
niebiesko.

% Lokalne style i makra do rysowania kolejnych kroków AVC na grafie instancji.
\tikzset{
	vcnode/.style={circle, draw, minimum size=0.55cm, inner sep=1pt, font=\small},
	cov/.style={fill=accent!25},
	sel/.style={accentB, line width=1.2pt},
	gone/.style={gray!45, dashed},
}
% #1..#6: dodatkowe style wierzchołków a,b,c,d,e,f (cov albo puste)
\newcommand{\vcnodes}[6]{%
	\node[vcnode,#1] (a) at (0,2) {$a$};
	\node[vcnode,#2] (b) at (2,2) {$b$};
	\node[vcnode,#3] (c) at (4,2) {$c$};
	\node[vcnode,#4] (d) at (0,0) {$d$};
	\node[vcnode,#5] (e) at (2,0) {$e$};
	\node[vcnode,#6] (f) at (4,0) {$f$};
}
% #1..#8: style krawędzi ab,bc,ad,be,cf,de,ef,bd (sel/gone albo puste = nadal w F)
\newcommand{\vcedges}[8]{%
	\draw[#1] (a)--(b); \draw[#2] (b)--(c); \draw[#3] (a)--(d);
	\draw[#4] (b)--(e); \draw[#5] (c)--(f); \draw[#6] (d)--(e);
	\draw[#7] (e)--(f); \draw[#8] (b)--(d);
}
\newcommand{\vcpanel}[1]{%
	\begin{tikzpicture}[scale=0.6, baseline=(current bounding box.center)]#1\end{tikzpicture}}
\newcommand{\vclabel}[1]{\node[font=\footnotesize] at (2,-0.9) {#1};}

\begin{dygresja}
	Przebieg pętli ($F$ -- niepokryte krawędzie, $B$ -- wybierane krawędzie):
	\begin{enumerate}[nosep]
		\item $F = \{ab, ad, bc, bd, be, cf, de, ef\}$. Wybór $ab$; usuwamy
		      $ab, ad, bc, bd, be$ (incydentne z~$a$ lub $b$).
		\item $F = \{cf, de, ef\}$. Wybór $cf$; usuwamy $cf, ef$ (incydentne z~$f$).
		\item $F = \{de\}$. Wybór $de$; usuwamy $de$ --- $F = \emptyset$.
	\end{enumerate}
\end{dygresja}

\begin{azfigure}
	\centering
	\vcpanel{\vcnodes{}{}{}{}{}{}\vcedges{}{}{}{}{}{}{}{}\vclabel{start}}%
	\;$\rightarrow$\;%
	\vcpanel{\vcnodes{cov}{cov}{}{}{}{}\vcedges{sel}{gone}{gone}{gone}{}{}{}{gone}\vclabel{wybór $ab$}}%
	\;$\rightarrow$\;%
	\vcpanel{\vcnodes{cov}{cov}{cov}{}{}{cov}\vcedges{sel}{gone}{gone}{gone}{sel}{}{gone}{gone}\vclabel{wybór $cf$}}%
	\;$\rightarrow$\;%
	\vcpanel{\vcnodes{cov}{cov}{cov}{cov}{cov}{cov}\vcedges{sel}{gone}{gone}{gone}{sel}{sel}{gone}{gone}\vclabel{wybór $de$}}
\end{azfigure}

\paragraph{Wynik.}
Algorytm dokłada oba końce trzech wybranych krawędzi $B = \{ab, cf, de\}$,
zwracając pokrycie
\[
	\boxed{C = \{a,\ b,\ c,\ d,\ e,\ f\}}
\]
o~liczności $|C| = 2|B| = 6$ --- czyli wszystkie wierzchołki grafu. Optymalne
pokrycie ma rozmiar $3$, np.\ $C^* = \{b, d, f\}$ (każda krawędź ma koniec
w~$b$, $d$ lub $f$). Ta instancja realizuje więc najgorszy przypadek: stosunek
$|C| / |C^*| = 2$, dokładnie gwarancja aproksymacji \textsc{AVC}.

% ── Rozwiązanie: Exact Subset Sum -- instancja 1 ─────────────────
\subsection{Exact Subset Sum -- instancja 1}

Dane: $S = \{1, 4, 5, 11, 12\}$ (w~kolejności $x_1, \ldots, x_5$) oraz $t = 21$.
Algorytm \textsc{ExactSubsetSum} startuje z~$L_0 = (0)$ i~w~$i$-tej iteracji
scala $L_{i-1}$ z~$L_{i-1} + x_i$ (procedura \textsc{MergeLists} -- posortowane
złączenie bez duplikatów), a~następnie usuwa z~wyniku wszystkie liczby $> t$.
Lista $L_i$ to posortowane sumy podzbiorów $\{x_1, \ldots, x_i\}$ nieprzekraczające
$t$. Element wprowadzony przez przesunięcie $L_{i-1} + x_i$ wyróżniono na~%
\textcolor{accentB}{pomarańczowo}.

\begin{dygresja}
	Przebieg iteracji ($L_{i-1} + x_i$ to składnik przesunięty):
	\begin{enumerate}[nosep]
		\item $x_1 = 1$:\quad
		      $(0) \cup (\textcolor{accentB}{1}) = (0, \textcolor{accentB}{1})$
		      \hfill $L_1 = (0, 1)$
		\item $x_2 = 4$:\quad
		      $(0, 1) \cup (\textcolor{accentB}{4}, \textcolor{accentB}{5})
		      = (0, 1, \textcolor{accentB}{4}, \textcolor{accentB}{5})$
		      \hfill $L_2 = (0, 1, 4, 5)$
		\item $x_3 = 5$:\quad
		      $(0, 1, 4, 5) \cup (\textcolor{accentB}{6}, \textcolor{accentB}{9},
		      \textcolor{accentB}{10})$ \\
		      (przesunięte $5, 6, 9, 10$; $5$ już było)
		      \hfill $L_3 = (0, 1, 4, 5, 6, 9, 10)$
		\item $x_4 = 11$:\quad scalenie z~$(11, 12, 15, 16, 17, 20, 21)$, nic $> 21$
		      \hfill\\
		      $L_4 = (0, 1, 4, 5, 6, 9, 10, \textcolor{accentB}{11},
		      \textcolor{accentB}{12}, \textcolor{accentB}{15},
		      \textcolor{accentB}{16}, \textcolor{accentB}{17},
		      \textcolor{accentB}{20}, \textcolor{accentB}{21})$
		\item $x_5 = 12$:\quad przesunięcie $L_4 + 12$ daje
		      $12, 13, 16, 17, 18, 21, 22, 23, 24, 27, 28, 29, 32, 33$;
		      po~usunięciu $> 21$ dochodzą tylko $\textcolor{accentB}{13}$
		      i~$\textcolor{accentB}{18}$ \\
		      $L_5 = (0, 1, 4, 5, 6, 9, 10, 11, 12, \textcolor{accentB}{13}, 15,
		      16, 17, \textcolor{accentB}{18}, 20, 21)$
	\end{enumerate}
\end{dygresja}

\paragraph{Wynik.}
Największy element listy $L_5$ nieprzekraczający $t = 21$ to
\[
	\boxed{z = 21}
\]
czyli dokładnie $t$ -- istnieje podzbiór o~sumie równej progowi, np.\
$\{4, 5, 12\}$ lub $\{1, 4, 5, 11\}$. Wersja decyzyjna ma więc odpowiedź
\emph{tak}, a~optymalna suma $\leq t$ wynosi $21$.

% ── Rozwiązanie: Approx Subset Sum -- instancja 1 ────────────────
\subsection{Approx Subset Sum -- instancja 1}

Dane: $S = \{106, 103, 205, 104\}$ (w~kolejności $x_1, \ldots, x_4$),
$t = 312$ oraz $\varepsilon = 0{,}4$. Wtedy $n = 4$ i~próg przycinania to
$\delta = \frac{\varepsilon}{2n} = \frac{0{,}4}{8} = 0{,}05$ (mnożnik
$1 + \delta = 1{,}05$). Startujemy z~$L_0 = (0)$.

W~każdej iteracji najpierw scalamy $L_{i-1}$ z~$L_{i-1} + x_i$, potem usuwamy
elementy $> t = 312$, a~na~końcu przycinamy. Przy przycinaniu zachowujemy
element $y$, gdy $y > last \cdot 1{,}05$, gdzie $last$ to ostatnio zachowany
element.
\begin{dygresja} \leavevmode
	\begin{itemize}[nosep]
		\item \textbf{$i = 1$, $x_1 = 106$.} \\
		      scalenie: $(0) \cup (106) = (0, 106)$; \\
		      po~usunięciu $> 312$: $(0, 106)$; \\
		      przycinanie: \\
		      \begin{tblr}{colspec={r c l l}, column{4}={font=\itshape}, rowsep=1pt}
			      $0$   &     &                     & zachowane, $last = 0$   \\
			      $106$ & $>$ & $0 \cdot 1{,}05 = 0$ & zachowane, $last = 106$ \\
		      \end{tblr} \\
		      $L_1 = (0, 106)$.
		\item \textbf{$i = 2$, $x_2 = 103$.} \\
		      scalenie: $(0, 106) \cup (103, 209) = (0, 103, 106, 209)$; \\
		      po~usunięciu $> 312$: bez zmian; \\
		      przycinanie: \\
		      \begin{tblr}{colspec={r c l l}, column{4}={font=\itshape}, rowsep=1pt}
			      $0$   &        &                              & zachowane, $last = 0$   \\
			      $103$ & $>$    & $0$                          & zachowane, $last = 103$ \\
			      $106$ & $\leq$ & $103 \cdot 1{,}05 = 108{,}15$ & usunięte                \\
			      $209$ & $>$    & $108{,}15$                   & zachowane, $last = 209$ \\
		      \end{tblr} \\
		      $L_2 = (0, 103, 209)$.
		\item \textbf{$i = 3$, $x_3 = 205$.} \\
		      scalenie: $(0, 103, 209) \cup (205, 308, 414) = (0, 103, 205, 209, 308, 414)$; \\
		      po~usunięciu $> 312$: $(0, 103, 205, 209, 308)$; \\
		      przycinanie: \\
		      \begin{tblr}{colspec={r c l l}, column{4}={font=\itshape}, rowsep=1pt}
			      $0$   &        &                              & zachowane, $last = 0$   \\
			      $103$ & $>$    & $0$                          & zachowane, $last = 103$ \\
			      $205$ & $>$    & $103 \cdot 1{,}05 = 108{,}15$ & zachowane, $last = 205$ \\
			      $209$ & $\leq$ & $205 \cdot 1{,}05 = 215{,}25$ & usunięte               \\
			      $308$ & $>$    & $215{,}25$                   & zachowane, $last = 308$ \\
		      \end{tblr} \\
		      $L_3 = (0, 103, 205, 308)$.
		\item \textbf{$i = 4$, $x_4 = 104$.} \\
		      scalenie: $(0, 103, 205, 308) \cup (104, 207, 309, 412)
		      = (0, 103, 104, 205, 207, 308, 309, 412)$; \\
		      po~usunięciu $> 312$: $(0, 103, 104, 205, 207, 308, 309)$; \\
		      przycinanie: \\
		      \begin{tblr}{colspec={r c l l}, column{4}={font=\itshape}, rowsep=1pt}
			      $0$   &        &                              & zachowane, $last = 0$   \\
			      $103$ & $>$    & $0$                          & zachowane, $last = 103$ \\
			      $104$ & $\leq$ & $103 \cdot 1{,}05 = 108{,}15$ & usunięte               \\
			      $205$ & $>$    & $108{,}15$                   & zachowane, $last = 205$ \\
			      $207$ & $\leq$ & $205 \cdot 1{,}05 = 215{,}25$ & usunięte               \\
			      $308$ & $>$    & $215{,}25$                   & zachowane, $last = 308$ \\
			      $309$ & $\leq$ & $308 \cdot 1{,}05 = 323{,}4$ & usunięte                \\
		      \end{tblr} \\
		      $L_4 = (0, 103, 205, 308)$.
	\end{itemize}
\end{dygresja}

\paragraph{Wynik.}
Zwracamy największy element listy $L_4$:
\[
	\boxed{z = 308}
\]
realizowany podzbiorem $\{103, 205\}$. Optimum w~granicy progu to
$z^* = 311$ (podzbiór $\{106, 205\}$): zostało utracone już w~iteracji
$i = 2$, gdy przycinanie usunęło $106$ jako bliskie zachowanemu $103$. Iloraz
$\frac{z^*}{z} = \frac{311}{308} \approx 1{,}010 \leq 1 + \varepsilon = 1{,}4$
-- zgodnie z~gwarancją aproksymacji.

% ── Rozwiązanie: Maximum Matching -- instancja 1 ─────────────────
\subsection{Maximum Matching -- instancja 1}

Graf dwudzielny $G = (V_1, V_2; E)$ z~$V_1 = \{a, b, c, d\}$, $V_2 = \{p, q, r, s\}$
oraz krawędziami $ap, aq, bp, br, cq, cr, cs, ds$. Algorytm \textsc{MM} startuje
z~$M = \emptyset$ i~dopóki \textsc{APS} zwraca ścieżkę powiększającą $P$,
podmienia $M \Assign M \div E(P)$ (różnica symetryczna -- krawędzie wolne na
$P$ wchodzą do~$M$, skojarzone wychodzą). \textsc{APS} szuka ścieżki w~grafie
skierowanym $G_M$ (krawędzie wolne kierowane $V_1 \to V_2$, skojarzone
$V_2 \to V_1$) z~wierzchołka wolnego w~$V_1$ do~wolnego w~$V_2$; przyjmujemy
przeszukiwanie w~głąb, odwiedzając sąsiadów alfabetycznie. Na~rysunkach
pierścień otacza wierzchołki \emph{wolne}, a~strzałki pokazują orientację
krawędzi w~$G_M$ (skojarzone w~górę, do~$V_1$; wolne w~dół, do~$V_2$).
Każdy wiersz pokazuje $G_M$ ze~znalezioną ścieżką powiększającą
(\textcolor{path}{na~różowo}, różowy pierścień to~wierzchołek startowy),
a~po~strzałce $\Rightarrow$ --- wynikowe skojarzenie $M$
(\textcolor{accentB}{na~pomarańczowo}).

% Lokalne makra do rysowania grafu instancji z~zaznaczonym skojarzeniem.
\tikzset{
	mmnode/.style={circle, fill=black, inner sep=1.5pt},
	sel/.style={accentB, line width=1.4pt, <-},  % skojarzona: strzałka V_2->V_1 (w górę)
	augF/.style={path, line width=1.4pt},        % na ścieżce, wolna (w dół, domyślne ->)
	augM/.style={path, line width=1.4pt, <-},    % na ścieżce, skojarzona (w górę)
}
% Strzałki w G_M: krawędź wolna ->, skojarzona przejmuje <- ze stylu sel.
% Pierścień wokół wierzchołków wolnych (lista w #1, pusta dozwolona).
\newcommand{\mmfreeR}[1]{\ifx\relax#1\relax\else
	\foreach \n in {#1} {\draw[black] (\n) circle (3.5pt);}\fi}
% #1..#8: style krawędzi ap,aq,bp,br,cq,cr,cs,ds (sel albo puste)
\newcommand{\mmedges}[8]{%
	\foreach \x/\n in {0/a, 1/b, 2/c, 3/d} {\node[mmnode, label=above:$\n$] (\n) at (\x,1) {};}
	\foreach \x/\n in {0/p, 1/q, 2/r, 3/s} {\node[mmnode, label=below:$\n$] (\n) at (\x,0) {};}
	\draw[->,#1] (a)--(p); \draw[->,#2] (a)--(q); \draw[->,#3] (b)--(p); \draw[->,#4] (b)--(r);
	\draw[->,#5] (c)--(q); \draw[->,#6] (c)--(r); \draw[->,#7] (c)--(s); \draw[->,#8] (d)--(s);}
% #1: style krawędzi, #2: etykieta, #3: lista wolnych, #4: początek ścieżki (pierścień różowy, pusty dozwolony)
\newcommand{\mmpanel}[4]{%
	\begin{tikzpicture}[x=0.8cm, y=0.9cm, scale=0.78, >=stealth, baseline=(current bounding box.center)]
		\mmedges#1\mmfreeR{#3}%
		\ifx\relax#4\relax\else\draw[path, line width=1.4pt] (#4) circle (3.5pt);\fi
		\node[font=\footnotesize] at (1.5,-0.95) {#2};\end{tikzpicture}}

\begin{dygresja}
	Przebieg pętli (ścieżki powiększające na~\textcolor{path}{różowo},
	krawędzie w~skojarzeniu na~\textcolor{accentB}{pomarańczowo}):
	\begin{enumerate}[nosep]
		\item $M = \emptyset$, wszystkie wierzchołki wolne. \textsc{APS} z~$a$
		      znajduje wolną krawędź $P_1 = a\,p$.
		      \hfill $M = \{ap\}$
		\item Wolne: $b, c, d$ oraz $q, r, s$. W~$G_M$ z~$b$: krawędź wolna
		      $b\!\to\!p$, dalej skojarzona $p\!\to\!a$, wolna $a\!\to\!q$ ($q$
		      wolne). Ścieżka powiększająca $P_2 = b\,p\,a\,q$ (zamiana: $ap$
		      wychodzi, $bp, aq$ wchodzą).
		      \hfill $M = \{aq, bp\}$
		\item Wolne: $c, d$ oraz $r, s$. Z~$c$: $c\!\to\!q$, $q\!\to\!a$,
		      $a\!\to\!p$, $p\!\to\!b$, $b\!\to\!r$ ($r$ wolne). Ścieżka
		      $P_3 = c\,q\,a\,p\,b\,r$ ($aq, bp$ wychodzą, $cq, ap, br$ wchodzą).
		      \hfill $M = \{ap, br, cq\}$
		\item Wolne: $d$ oraz $s$. Jedyna krawędź $d$ to wolna $d\,s$:
		      $P_4 = d\,s$.
		      \hfill $M = \{ap, br, cq, ds\}$
		\item Brak wolnych wierzchołków w~$V_1$ -- \textsc{APS} zwraca NULL,
		      pętla się kończy.
	\end{enumerate}
\end{dygresja}

% Każdy wiersz: G_M ze ścieżką powiększającą (\textcolor{path}{na różowo}, pierścień
% startowy też różowy) -> wynikowe M. Kolumny w tblr wyrównują rysunki w pionie.
\begin{azfigure}
	\centering
	\begin{tblr}{colspec={ccc}, rowsep=10pt, colsep=6pt}
		\mmpanel{{augF}{}{}{}{}{}{}{}}{$P_1 = a\,p$}{a,b,c,d,p,q,r,s}{a}
		& $\Rightarrow$ &
		\mmpanel{{sel}{}{}{}{}{}{}{}}{$M = \{ap\}$}{b,c,d,q,r,s}{} \\
		\mmpanel{{augM}{augF}{augF}{}{}{}{}{}}{$P_2 = b\,p\,a\,q$}{b,c,d,q,r,s}{b}
		& $\Rightarrow$ &
		\mmpanel{{}{sel}{sel}{}{}{}{}{}}{$M = \{aq, bp\}$}{c,d,r,s}{} \\
		\mmpanel{{augF}{augM}{augM}{augF}{augF}{}{}{}}{$P_3 = c\,q\,a\,p\,b\,r$}{c,d,r,s}{c}
		& $\Rightarrow$ &
		\mmpanel{{sel}{}{}{sel}{sel}{}{}{}}{$M = \{ap, br, cq\}$}{d,s}{} \\
		\mmpanel{{sel}{}{}{sel}{sel}{}{}{augF}}{$P_4 = d\,s$}{d,s}{d}
		& $\Rightarrow$ &
		\mmpanel{{sel}{}{}{sel}{sel}{}{}{sel}}{$M = \{ap, br, cq, ds\}$}{}{} \\
	\end{tblr}
\end{azfigure}

\paragraph{Wynik.}
Algorytm zwraca skojarzenie
\[
	\boxed{M = \{ap,\ br,\ cq,\ ds\}}
\]
o~liczności $|M| = 4$. Jest ono \emph{doskonałe} -- każdy wierzchołek jest
skojarzony -- więc i~najliczniejsze (większe niż $|V_1| = 4$ być nie może).
Dwie środkowe iteracje pokazują istotę metody: zachłanne krawędzie $ap$, a~potem
$aq, bp$ blokowałyby kolejne wierzchołki, ale ścieżki powiększające $P_2$ i~$P_3$
przekładają skojarzenie wzdłuż naprzemiennych krawędzi i~za~każdym razem
zwiększają $|M|$ o~$1$.

% ── Rozwiązanie: Para najbliższych punktów -- instancja 2 ──────────
\subsection{Para najbliższych punktów -- instancja 2}

% Kolory grup w drzewie rekursji (4 liście = 4 wywołania siłowe).
\colorlet{grpLLb}{accent}%  {p1,p2}
\colorlet{grpLRb}{accentC}% {p3,p4}
\colorlet{grpRLb}{accentB}% {p5,p6}
\colorlet{grpRRb}{violet}%  {p7,p8}
\colorlet{grpLb}{accent}%   lewa połowa P_L
\colorlet{grpRb}{accentB}%  prawa połowa P_R

Dane (instancja podana tabelą), $n = 8$. Ponumerujmy punkty w~kolejności
wejściowej:
\[
	p_1=(0,0),\ p_2=(1,6),\ p_3=(4,2),\ p_4=(5,5),\
	p_5=(7,1),\ p_6=(8,7),\ p_7=(9,3),\ p_8=(11,6).
\]

\paragraph{Sortowanie wstępne.}
Tablice $X$ (rosnąco po~$x$) i~$Y$ (rosnąco po~$y$):
\begin{align*}
	X &= \big[\,p_1(0,0),\ p_2(1,6),\ p_3(4,2),\ p_4(5,5),\ p_5(7,1),\ p_6(8,7),\ p_7(9,3),\ p_8(11,6)\,\big], \\
	Y &= \big[\,p_1(0,0),\ p_5(7,1),\ p_3(4,2),\ p_7(9,3),\ p_4(5,5),\ p_2(1,6),\ p_8(11,6),\ p_6(8,7)\,\big].
\end{align*}

\paragraph{Dziel.}
$|P|=8>3$, dzielimy prostą $l\colon x=6$ (mediana współrzędnej~$x$:
$\frac{X(4).x + X(5).x}{2} = \frac{5+7}{2} = 6$; zob.~uwagę o~wyborze~$l$)
na~$|P_L|=|P_R|=4$, a~każdą połowę dalej,
aż do podzbiorów po~$2$ punkty (liście --- sprawdzane siłowo). Strukturę rekursji
zbiera rysunek~\ref{fig:closest2:tree}.

\begin{azfigure}
	\centering
	\begin{tikzpicture}[
			font=\small,
			level distance=12mm,
			level 1/.style={sibling distance=42mm},
			level 2/.style={sibling distance=22mm},
			grp/.style={draw, rounded corners=2pt, inner sep=3pt, thick},
			edge from parent/.style={draw, gray, -}]
		\node[grp] {$P$ \scriptsize(8)}
			child {node[grp, grpLb, text=grpLb] {$P_L$ \scriptsize(4)}
				child {node[grp, grpLLb, text=grpLLb] {$\{p_1,p_2\}$}}
				child {node[grp, grpLRb, text=grpLRb] {$\{p_3,p_4\}$}}}
			child {node[grp, grpRb, text=grpRb] {$P_R$ \scriptsize(4)}
				child {node[grp, grpRLb, text=grpRLb] {$\{p_5,p_6\}$}}
				child {node[grp, grpRRb, text=grpRRb] {$\{p_7,p_8\}$}}};
	\end{tikzpicture}
	\caption{Drzewo rekursji. Cztery liście (po~2 punkty) liczone są siłowo;
		ich kolory powtarzają się na~rysunku~\ref{fig:closest2:result}.}
	\label{fig:closest2:tree}
\end{azfigure}

\[
	P_L=\{\textcolor{grpLLb}{p_1,p_2},\ \textcolor{grpLRb}{p_3,p_4}\},\qquad
	P_R=\{\textcolor{grpRLb}{p_5,p_6},\ \textcolor{grpRRb}{p_7,p_8}\}.
\]

\paragraph{Zwyciężaj --- lewa połowa $P_L$.}
Liście (po~2 punkty --- siłowo) i~faza \emph{połącz} wzdłuż prostej $x=2{,}5$:
\[
	d(\textcolor{grpLLb}{p_1,p_2})=\sqrt{1^2+6^2}=\sqrt{37}\approx6{,}08,\qquad
	\boxed{d(\textcolor{grpLRb}{p_3,p_4})=\sqrt{1^2+3^2}=\sqrt{10}\approx3{,}16.}
\]
Pas o~szerokości $2\sqrt{10}$ obejmuje wszystkie cztery punkty, ale żadna para
„międzystronowa” nie jest bliższa niż $\sqrt{10}$. Zatem $\delta_L=\sqrt{10}$ dla
pary $\big(p_3(4,2),\,p_4(5,5)\big)$.

\paragraph{Zwyciężaj --- prawa połowa $P_R$.}
Liście i~faza \emph{połącz} wzdłuż $x=8{,}5$:
\[
	d(\textcolor{grpRLb}{p_5,p_6})=\sqrt{1^2+6^2}=\sqrt{37}\approx6{,}08,\qquad
	d(\textcolor{grpRRb}{p_7,p_8})=\sqrt{2^2+3^2}=\sqrt{13}\approx3{,}61.
\]
W~pasie wokół $x=8{,}5$ leżą $p_5,p_6,p_7,p_8$; sprawdzając pary stwierdzamy
\[
	\boxed{d(p_5,p_7)=\sqrt{2^2+2^2}=\sqrt{8}=2\sqrt2\approx2{,}83,}
\]
co poprawia $\delta$. Zatem $\delta_R=\sqrt8$ dla pary $\big(p_5(7,1),\,p_7(9,3)\big)$.

\paragraph{Połącz --- poziom górny.}
$\delta=\min(\delta_L,\delta_R)=\min(\sqrt{10},\sqrt8)=\sqrt8$. Pas wokół
$l\colon x=6$ ma szerokość $2\sqrt8\approx5{,}66$: $x\in[\,6-2\sqrt2,\ 6+2\sqrt2\,]
\approx[3{,}17,\ 8{,}83]$. W~pasie (rosnąco po~$y$) leżą $p_5(7,1),p_3(4,2),
p_4(5,5),p_6(8,7)$; sąsiednie pary mają odległości $\geq\sqrt{10}>\sqrt8$, więc
$\delta$ bez zmian.

\paragraph{Wynik.}
Para najbliższych punktów to $(7,1)$ i~$(9,3)$ w~odległości
$\delta=2\sqrt2\approx2{,}83$.

\begin{azfigure}
	\centering
	\begin{tikzpicture}[>=stealth, font=\small, scale=0.45]
		\draw[gray!50, very thin] (0,0) grid (11,7);
		\draw[->] (0,0)--(12,0) node[right]{$x$};
		\draw[->] (0,0)--(0,8) node[above]{$y$};
		\draw[thick, accent] (6,-0.5)--(6,7.5) node[above]{$l$};
		\foreach \x/\y/\c in {0/0/grpLLb,1/6/grpLLb,4/2/grpLRb,5/5/grpLRb,%
				7/1/grpRLb,8/7/grpRLb,9/3/grpRRb,11/6/grpRRb}
			\fill[\c] (\x,\y) circle (3.5pt);
		\draw[red, thick] (7,1)--(9,3);
		\draw[red, thick] (7,1) circle (6pt);
		\draw[red, thick] (9,3) circle (6pt);
		\node[red, below=3pt] at (7,1) {$(7,1)$};
		\node[red, right=3pt] at (9,3) {$(9,3)$};
	\end{tikzpicture}
	\caption{Instancja~2: prosta podziału $l\colon x=6$, punkty pokolorowane wg
		liścia drzewa rekursji oraz znaleziona para $(7,1)$, $(9,3)$ (czerwona).}
	\label{fig:closest2:result}
\end{azfigure}

% ── Rozwiązanie: Matrix-Chain-Order -- instancja 2 ────────────────
\subsection{Matrix-Chain-Order -- instancja 2}

Dane: $n=5$, $p=(10,\ 3,\ 12,\ 5,\ 50,\ 6)$, czyli
\[
	A_1\ (10\times 3),\quad A_2\ (3\times 12),\quad A_3\ (12\times 5),\quad
	A_4\ (5\times 50),\quad A_5\ (50\times 6).
\]
Wypełniamy tablice $m$ i~$s$ \emph{przekątnymi} po~rosnącej długości podciągu
$L$; w~każdym kroku dopisane komórki wyróżniono na
\textcolor{accentB}{pomarańczowo}, a~$\times$ to pola nieużywane. Skrót:
$A_i=p_{i-1}\times p_i$.

\medskip
\noindent Stan początkowy ($L=1$, koszt $0$):\\[0.7em]
\begin{minipage}{0.45\linewidth}
	$m$
	$\begin{bNiceArray}{c|ccccc}[margin,hvlines]
			i \backslash j & 1 & 2 & 3 & 4 & 5 \\
			1 & 0 &  &  &  &  \\
			2 & \times & 0 &  &  &  \\
			3 & \times & \times & 0 &  &  \\
			4 & \times & \times & \times & 0 &  \\
			5 & \times & \times & \times & \times & 0
		\end{bNiceArray}$
\end{minipage}\hfill
\begin{minipage}{0.45\linewidth}
	$s$
	$\begin{bNiceArray}{c|ccccc}[margin,hvlines]
			i \backslash j & 1 & 2 & 3 & 4 & 5 \\
			1 & \times &  &  &  &  \\
			2 & \times & \times &  &  &  \\
			3 & \times & \times & \times &  &  \\
			4 & \times & \times & \times & \times &  \\
			5 & \times & \times & \times & \times & \times
		\end{bNiceArray}$
\end{minipage}

\medskip
\noindent{\color{placeholdercolor}\rule{\linewidth}{0.4pt}}\par\smallskip
\noindent{\color{dygresjagray}\footnotesize$\begin{array}{c|*{6}{c}} i & 0 & 1 & 2 & 3 & 4 & 5 \\ \hline p_i & 10 & 3 & 12 & 5 & 50 & 6 \end{array}$\qquad $A_i = p_{i-1}\times p_i$}
\begin{flalign*}
	L=2: \quad & m(1, 2) = 10 \cdot 3 \cdot 12 = 360  &  & \\
	           & m(2, 3) = 3 \cdot 12 \cdot 5 = 180   &  & \\
	           & m(3, 4) = 12 \cdot 5 \cdot 50 = 3000 &  & \\
	           & m(4, 5) = 5 \cdot 50 \cdot 6 = 1500  &  &
\end{flalign*}

\medskip
\noindent po $L=2$:\\[0.7em]
\begin{minipage}{0.45\linewidth}
	$m$
	$\begin{bNiceArray}{c|ccccc}[margin,hvlines]
			i \backslash j & 1 & 2 & 3 & 4 & 5 \\
			1 & 0 & \color{accentB}360 &  &  &  \\
			2 & \times & 0 & \color{accentB}180 &  &  \\
			3 & \times & \times & 0 & \color{accentB}3000 &  \\
			4 & \times & \times & \times & 0 & \color{accentB}1500 \\
			5 & \times & \times & \times & \times & 0
		\end{bNiceArray}$
\end{minipage}\hfill
\begin{minipage}{0.45\linewidth}
	$s$
	$\begin{bNiceArray}{c|ccccc}[margin,hvlines]
			i \backslash j & 1 & 2 & 3 & 4 & 5 \\
			1 & \times & \color{accentB}1 &  &  &  \\
			2 & \times & \times & \color{accentB}2 &  &  \\
			3 & \times & \times & \times & \color{accentB}3 &  \\
			4 & \times & \times & \times & \times & \color{accentB}4 \\
			5 & \times & \times & \times & \times & \times
		\end{bNiceArray}$
\end{minipage}

\medskip
\noindent{\color{placeholdercolor}\rule{\linewidth}{0.4pt}}\par\smallskip
\noindent{\color{dygresjagray}\footnotesize$\begin{array}{c|*{6}{c}} i & 0 & 1 & 2 & 3 & 4 & 5 \\ \hline p_i & 10 & 3 & 12 & 5 & 50 & 6 \end{array}$\qquad $A_i = p_{i-1}\times p_i$}
\begin{flalign*}
	L=3: \quad & m(1, 3) = \min \begin{cases}
		                            m(1,1) + m(2,3) + 10 \cdot 3 \cdot 5 \\
		                            m(1,2) + m(3,3) + 10 \cdot 12 \cdot 5
	                            \end{cases} = \min \begin{cases} 0+180+150 \\ 360+0+600 \end{cases} = \min \begin{cases} 330 \\ 960 \end{cases} = 330 \\
	           & m(2, 4) = \min \begin{cases}
		                            m(2,2) + m(3,4) + 3 \cdot 12 \cdot 50 \\
		                            m(2,3) + m(4,4) + 3 \cdot 5 \cdot 50
	                            \end{cases} = \min \begin{cases} 0+3000+1800 \\ 180+0+750 \end{cases} = \min \begin{cases} 4800 \\ 930 \end{cases} = 930 \\
	           & m(3, 5) = \min \begin{cases}
		                            m(3,3) + m(4,5) + 12 \cdot 5 \cdot 6 \\
		                            m(3,4) + m(5,5) + 12 \cdot 50 \cdot 6
	                            \end{cases} = \min \begin{cases} 0+1500+360 \\ 3000+0+3600 \end{cases} = \min \begin{cases} 1860 \\ 6600 \end{cases} = 1860
\end{flalign*}

\medskip
\noindent po $L=3$:\\[0.7em]
\begin{minipage}{0.45\linewidth}
	$m$
	$\begin{bNiceArray}{c|ccccc}[margin,hvlines]
			i \backslash j & 1 & 2 & 3 & 4 & 5 \\
			1 & 0 & 360 & \color{accentB}330 &  &  \\
			2 & \times & 0 & 180 & \color{accentB}930 &  \\
			3 & \times & \times & 0 & 3000 & \color{accentB}1860 \\
			4 & \times & \times & \times & 0 & 1500 \\
			5 & \times & \times & \times & \times & 0
		\end{bNiceArray}$
\end{minipage}\hfill
\begin{minipage}{0.45\linewidth}
	$s$
	$\begin{bNiceArray}{c|ccccc}[margin,hvlines]
			i \backslash j & 1 & 2 & 3 & 4 & 5 \\
			1 & \times & 1 & \color{accentB}1 &  &  \\
			2 & \times & \times & 2 & \color{accentB}3 &  \\
			3 & \times & \times & \times & 3 & \color{accentB}3 \\
			4 & \times & \times & \times & \times & 4 \\
			5 & \times & \times & \times & \times & \times
		\end{bNiceArray}$
\end{minipage}

\medskip
\noindent{\color{placeholdercolor}\rule{\linewidth}{0.4pt}}\par\smallskip
\noindent{\color{dygresjagray}\footnotesize$\begin{array}{c|*{6}{c}} i & 0 & 1 & 2 & 3 & 4 & 5 \\ \hline p_i & 10 & 3 & 12 & 5 & 50 & 6 \end{array}$\qquad $A_i = p_{i-1}\times p_i$}
\begin{flalign*}
	L=4: \quad & m(1, 4) = \min \begin{cases}
		                            m(1,1) + m(2,4) + 10 \cdot 3 \cdot 50 \\
		                            m(1,2) + m(3,4) + 10 \cdot 12 \cdot 50 \\
		                            m(1,3) + m(4,4) + 10 \cdot 5 \cdot 50
	                            \end{cases} = \min \begin{cases} 0+930+1500 \\ 360+3000+6000 \\ 330+0+2500 \end{cases} = \min \begin{cases} 2430 \\ 9360 \\ 2830 \end{cases} = 2430 \\
	           & m(2, 5) = \min \begin{cases}
		                            m(2,2) + m(3,5) + 3 \cdot 12 \cdot 6 \\
		                            m(2,3) + m(4,5) + 3 \cdot 5 \cdot 6 \\
		                            m(2,4) + m(5,5) + 3 \cdot 50 \cdot 6
	                            \end{cases} = \min \begin{cases} 0+1860+216 \\ 180+1500+90 \\ 930+0+900 \end{cases} = \min \begin{cases} 2076 \\ 1770 \\ 1830 \end{cases} = 1770
\end{flalign*}

\medskip
\noindent po $L=4$:\\[0.7em]
\begin{minipage}{0.45\linewidth}
	$m$
	$\begin{bNiceArray}{c|ccccc}[margin,hvlines]
			i \backslash j & 1 & 2 & 3 & 4 & 5 \\
			1 & 0 & 360 & 330 & \color{accentB}2430 &  \\
			2 & \times & 0 & 180 & 930 & \color{accentB}1770 \\
			3 & \times & \times & 0 & 3000 & 1860 \\
			4 & \times & \times & \times & 0 & 1500 \\
			5 & \times & \times & \times & \times & 0
		\end{bNiceArray}$
\end{minipage}\hfill
\begin{minipage}{0.45\linewidth}
	$s$
	$\begin{bNiceArray}{c|ccccc}[margin,hvlines]
			i \backslash j & 1 & 2 & 3 & 4 & 5 \\
			1 & \times & 1 & 1 & \color{accentB}1 &  \\
			2 & \times & \times & 2 & 3 & \color{accentB}3 \\
			3 & \times & \times & \times & 3 & 3 \\
			4 & \times & \times & \times & \times & 4 \\
			5 & \times & \times & \times & \times & \times
		\end{bNiceArray}$
\end{minipage}

\medskip
\noindent{\color{placeholdercolor}\rule{\linewidth}{0.4pt}}\par\smallskip
\noindent{\color{dygresjagray}\footnotesize$\begin{array}{c|*{6}{c}} i & 0 & 1 & 2 & 3 & 4 & 5 \\ \hline p_i & 10 & 3 & 12 & 5 & 50 & 6 \end{array}$\qquad $A_i = p_{i-1}\times p_i$}
\begin{flalign*}
	L=5: \quad & m(1, 5) = \min \begin{cases}
		                            m(1,1) + m(2,5) + 10 \cdot 3 \cdot 6 \\
		                            m(1,2) + m(3,5) + 10 \cdot 12 \cdot 6 \\
		                            m(1,3) + m(4,5) + 10 \cdot 5 \cdot 6 \\
		                            m(1,4) + m(5,5) + 10 \cdot 50 \cdot 6
	                            \end{cases} = \min \begin{cases} 0+1770+180 \\ 360+1860+720 \\ 330+1500+300 \\ 2430+0+3000 \end{cases} = \min \begin{cases} 1950 \\ 2940 \\ 2130 \\ 5430 \end{cases} = 1950
\end{flalign*}

\medskip
\noindent tablice końcowe:\\[0.7em]
\begin{minipage}{0.45\linewidth}
	$m$
	$\begin{bNiceArray}{c|ccccc}[margin,hvlines]
			i \backslash j & 1 & 2 & 3 & 4 & 5 \\
			1 & 0 & 360 & 330 & 2430 & \color{accentB}1950 \\
			2 & \times & 0 & 180 & 930 & 1770 \\
			3 & \times & \times & 0 & 3000 & 1860 \\
			4 & \times & \times & \times & 0 & 1500 \\
			5 & \times & \times & \times & \times & 0
		\end{bNiceArray}$
\end{minipage}\hfill
\begin{minipage}{0.45\linewidth}
	$s$
	$\begin{bNiceArray}{c|ccccc}[margin,hvlines]
			i \backslash j & 1 & 2 & 3 & 4 & 5 \\
			1 & \times & 1 & 1 & 1 & \color{accentB}1 \\
			2 & \times & \times & 2 & 3 & 3 \\
			3 & \times & \times & \times & 3 & 3 \\
			4 & \times & \times & \times & \times & 4 \\
			5 & \times & \times & \times & \times & \times
		\end{bNiceArray}$
\end{minipage}

\paragraph{Optymalne nawiasowanie.}
Odtwarzamy je z~tablicy $s$:
\[
	s(1,5)=1 \Rightarrow (A_1)(A_2 A_3 A_4 A_5),\quad
	s(2,5)=3 \Rightarrow (A_2 A_3)(A_4 A_5),\quad
	s(2,3)=2,\ s(4,5)=4.
\]
Stąd optymalne nawiasowanie
\[
	\boxed{A_1\,\big((A_2\,A_3)\,(A_4\,A_5)\big)}
\]
o~minimalnym koszcie $m(1,5)=1950$ mnożeń skalarnych.

% ── Rozwiązanie: LCS -- instancja 2 ───────────────────────────────
\subsection{Najdłuższy wspólny podciąg (LCS) -- instancja 2}

Dane: $X = (G,\ A,\ T,\ T,\ A,\ C,\ A)$, $Y = (G,\ C,\ A,\ T,\ A,\ A,\ T)$, czyli
$m = 7$, $n = 7$. Tablicę kosztów $c$ wypełniamy wierszami regułą jak w~instancji~1,
a~równolegle w~tablicy $b$ zapisujemy kierunek ($\nwarrow$ -- dopasowanie,
$\uparrow$ -- góra, $\leftarrow$ -- lewo; remisy na korzyść góry). Komórki
wchodzące do wybranego LCS wyróżniono na~\textcolor{accentB}{pomarańczowo}.

\begin{aztable}
	\centering
	\renewcommand{\arraystretch}{1.2}
	\begin{minipage}{0.48\linewidth}
		\centering
		Tablica kosztów $c$: \\[0.5em]
		\begin{NiceTabular}{cc|c|ccccccc}[margin,hvlines]
			    & $j$   & 0     & 1   & 2   & 3   & 4   & 5   & 6   & 7   \\
			$i$ &       & $y_j$ & $G$ & $C$ & $A$ & $T$ & $A$ & $A$ & $T$ \\
			\hline
			0   & $x_i$ & 0     & 0   & 0   & 0   & 0   & 0   & 0   & 0   \\
			1   & $G$   & 0     & \color{accentB}1 & 1 & 1 & 1 & 1 & 1 & 1 \\
			2   & $A$   & 0     & 1   & 1   & \color{accentB}2 & 2 & 2 & 2 & 2 \\
			3   & $T$   & 0     & 1   & 1   & 2   & 3   & 3   & 3   & 3   \\
			4   & $T$   & 0     & 1   & 1   & 2   & \color{accentB}3 & 3 & 3 & 4 \\
			5   & $A$   & 0     & 1   & 1   & 2   & 3   & \color{accentB}4 & 4 & 4 \\
			6   & $C$   & 0     & 1   & 2   & 2   & 3   & 4   & 4   & 4   \\
			7   & $A$   & 0     & 1   & 2   & 3   & 3   & 4   & \color{accentB}5 & 5 \\
		\end{NiceTabular}
	\end{minipage}\hfill
	\begin{minipage}{0.48\linewidth}
		\centering
		Tablica strzałek $b$: \\[0.5em]
		\begin{NiceTabular}{cc|c|ccccccc}[margin,hvlines]
			    & $j$   & 0     & 1          & 2            & 3            & 4            & 5            & 6            & 7            \\
			$i$ &       & $y_j$ & $G$        & $C$          & $A$          & $T$          & $A$          & $A$          & $T$          \\
			\hline
			0   & $x_i$ &       &            &              &              &              &              &              &              \\
			1   & $G$   &       & \color{accentB}$\nwarrow$ & $\leftarrow$ & $\leftarrow$ & $\leftarrow$ & $\leftarrow$ & $\leftarrow$ & $\leftarrow$ \\
			2   & $A$   &       & $\uparrow$ & $\uparrow$   & \color{accentB}$\nwarrow$ & $\leftarrow$ & $\nwarrow$   & $\nwarrow$   & $\leftarrow$ \\
			3   & $T$   &       & $\uparrow$ & $\uparrow$   & $\uparrow$   & $\nwarrow$   & $\leftarrow$ & $\leftarrow$ & $\nwarrow$   \\
			4   & $T$   &       & $\uparrow$ & $\uparrow$   & $\uparrow$   & \color{accentB}$\nwarrow$ & $\uparrow$ & $\uparrow$ & $\nwarrow$ \\
			5   & $A$   &       & $\uparrow$ & $\uparrow$   & $\nwarrow$   & $\uparrow$   & \color{accentB}$\nwarrow$ & $\nwarrow$ & $\uparrow$ \\
			6   & $C$   &       & $\uparrow$ & $\nwarrow$   & $\uparrow$   & $\uparrow$   & $\uparrow$   & $\uparrow$   & $\uparrow$   \\
			7   & $A$   &       & $\uparrow$ & $\uparrow$   & $\nwarrow$   & $\uparrow$   & $\nwarrow$   & \color{accentB}$\nwarrow$ & $\leftarrow$ \\
		\end{NiceTabular}
	\end{minipage}
\end{aztable}

\begin{dygresja}
	Po jednej komórce na każdą gałąź reguły,
	$X = (x_1, \dots, x_7)$, $Y = (y_1, \dots, y_7)$:
	\begin{itemize}
		\item \textbf{Dopasowanie} ($\nwarrow$): $c(1, 1)$, $x_1 = G = y_1$, więc
		      $c(1, 1) = c(0, 0) + 1 = 1$.
		\item \textbf{Góra} ($\uparrow$): $c(2, 1)$, $x_2 = A \neq G = y_1$; sąsiedzi
		      $c(1, 1) = 1 \geq c(2, 0) = 0$, więc $c(2, 1) = 1$.
		\item \textbf{Lewo} ($\leftarrow$): $c(1, 2)$, $x_1 = G \neq C = y_2$; sąsiedzi
		      $c(0, 2) = 0 < c(1, 1) = 1$, więc $c(1, 2) = 1$.
		\item \textbf{Remis}: $c(2, 7)$, $A \neq T$, $c(1, 7) = c(2, 6) = 2$ -- warunek
		      $\geq$ spełniony, więc $\uparrow$.
	\end{itemize}
\end{dygresja}

\paragraph{Odtworzenie podciągu.}
Startujemy z~$(m, n) = (7, 7)$ i~idziemy za strzałkami; przy $\nwarrow$ schodzimy
po~przekątnej, dopisując $x_i$ do~LCS:
\[
	\begin{aligned}
		(7, 7) & \xrightarrow{\leftarrow} (7, 6) \xrightarrow{\nwarrow} (6, 5) \;[\,x_7 = A\,]
		\xrightarrow{\uparrow} (5, 5) \xrightarrow{\nwarrow} (4, 4) \;[\,x_5 = A\,] \\
		       & \xrightarrow{\nwarrow} (3, 3) \;[\,x_4 = T\,]
		\xrightarrow{\uparrow} (2, 3) \xrightarrow{\nwarrow} (1, 2) \;[\,x_2 = A\,]
		\xrightarrow{\leftarrow} (1, 1) \xrightarrow{\nwarrow} (0, 0) \;[\,x_1 = G\,].
	\end{aligned}
\]
Znaki dopisywane są \emph{po} powrocie z~rekurencji, więc w~kolejności od
najgłębszego: $G, A, T, A, A$. Stąd najdłuższy wspólny podciąg
\[
	\boxed{\text{LCS}(X, Y) = (G,\ A,\ T,\ A,\ A)}
\]
o~długości $c(7, 7) = 5$.

% ── Rozwiązanie: Zachłanny wybór zajęć -- instancja 2 ─────────────
\subsection{Zachłanny wybór zajęć -- instancja 2}

Dane ($n = 9$ zajęć). Tym razem nie są one posortowane po~czasie zakończenia:
$t_2 = 5 > t_3 = 4$. Najpierw więc sortujemy rosnąco po~$t_i$ (miejscami zamieniają się
$z_2$ i~$z_3$ oraz $z_4$ i~$z_5$):

\begin{aztable}
	\centering
	\renewcommand{\arraystretch}{1.2}
	\begin{NiceTabular}{c|ccccccccc}[margin,hvlines]
		$i$   & \color{accentB}1 & 3 & 2 & \color{accentB}5 & 4 & \color{accentB}6 & 7 & \color{accentB}8 & 9  \\
		\hline
		$s_i$ & \color{accentB}0 & 1 & 2 & \color{accentB}4 & 6 & \color{accentB}8 & 7 & \color{accentB}9 & 11 \\
		$t_i$ & \color{accentB}3 & 4 & 5 & \color{accentB}7 & 8 & \color{accentB}9 & 11& \color{accentB}12& 14
	\end{NiceTabular}
\end{aztable}

Algorytm \textsc{GreedyActivitySelector} wybiera pierwsze zajęcie i~idzie dalej,
dodając $z_i$ wtedy i~tylko wtedy, gdy $s_i \geq t_j$ dla ostatnio wybranego
$z_j$. Wybrane zajęcia wyróżniono na~\textcolor{accentB}{pomarańczowo}.

\begin{dygresja}
	Przebieg pętli (kolejność po~sortowaniu: $z_1, z_3, z_2, z_5, z_4, z_6, z_7, z_8, z_9$;
	początkowo $A = \{z_1\}$, $j = 1$, $t_j = 3$):
	\begin{itemize}
		\item $z_3$: $s_3 = 1 \geq t_j = 3$? \emph{nie} -- pomijamy.
		\item $z_2$: $s_2 = 2 \geq 3$? \emph{nie} -- pomijamy.
		\item $z_5$: $s_5 = 4 \geq 3$? \emph{tak} -- $A \cup \{z_5\}$, $j \Assign 5$, $t_j = 7$.
		\item $z_4$: $s_4 = 6 \geq 7$? \emph{nie} -- pomijamy.
		\item $z_6$: $s_6 = 8 \geq 7$? \emph{tak} -- $A \cup \{z_6\}$, $j \Assign 6$, $t_j = 9$.
		\item $z_7$: $s_7 = 7 \geq 9$? \emph{nie} -- pomijamy.
		\item $z_8$: $s_8 = 9 \geq 9$? \emph{tak} -- $A \cup \{z_8\}$, $j \Assign 8$, $t_j = 12$.
		\item $z_9$: $s_9 = 11 \geq 12$? \emph{nie} -- pomijamy.
	\end{itemize}
\end{dygresja}

\paragraph{Wynik.}
Wybrane zajęcia są parami zgodne, bo ich przedziały
$\langle 0, 3)$, $\langle 4, 7)$, $\langle 8, 9)$, $\langle 9, 12)$ są rozłączne:
\[
	\boxed{A = \{z_1,\ z_5,\ z_6,\ z_8\}}
\]
o~liczności $|A| = 4$.

% ── Rozwiązanie: Huffman -- instancja 2 ──────────────────────────
\subsection{Algorytm Huffmana -- instancja 2}

Dane: kolejka startowa $(g, 2),\ (o, 3),\ (k, 6),\ (m, 8),\ (l, 11),\ (u, 14),\ (e, 20)$,
sumaryczna częstość $64$. W~każdej z~$n - 1 = 6$ iteracji łączymy dwa korzenie
o~\emph{najmniejszej} częstości (mniejszy jako lewe dziecko, krawędź $0$; remisy
wg~kolejności w~kolejce), wstawiając nowy wierzchołek z~powrotem. Style drzew
i~makro \texttt{\textbackslash huffstep} przejęte z~instancji~1.

\tikzset{
	htree/.style={
			level distance=0.95cm,
			level 1/.style={sibling distance=3.4cm},
			level 2/.style={sibling distance=1.9cm},
			level 3/.style={sibling distance=1.2cm},
			level 4/.style={sibling distance=1.0cm},
			every node/.style={circle, draw, minimum size=0.7cm, inner sep=0pt, font=\footnotesize, align=center},
			elabel left/.style={draw=none, font=\scriptsize, inner sep=1pt, midway, sloped, above},
			elabel right/.style={draw=none, font=\scriptsize, inner sep=1pt, midway, sloped, above},
		}
}

\huffstep{%
	\textbf{1.} Łączymy $g$ i~$o$. \\
	Zostają: $(5), (k, 6), (m, 8), (l, 11), (u, 14), (e, 20)$.
}{%
	\begin{tikzpicture}[htree]
		\node {5}
		child { node {g,2} }
		child { node {o,3} };
	\end{tikzpicture}
}

\huffstep{%
	\textbf{2.} Łączymy $5$ i~$k$. \\
	Zostają: $(m, 8), (11), (l, 11), (u, 14), (e, 20)$.
}{%
	\begin{tikzpicture}[htree]
		\node {11}
		child { node {5}
				child { node {g,2} }
				child { node {o,3} }
			}
		child { node {k,6} };
	\end{tikzpicture}
}

\huffstep{%
	\textbf{3.} Łączymy $m$ i~$l$. \\
	Zostają: $(11), (u, 14), (19), (e, 20)$.
}{%
	\begin{tikzpicture}[htree]
		\node {19}
		child { node {m,8} }
		child { node {l,11} };
	\end{tikzpicture}
}

\huffstep{%
	\textbf{4.} Łączymy $11$ i~$u$. \\
	Zostają: $(19), (e, 20), (25)$.
}{%
	\begin{tikzpicture}[htree]
		\node {25}
		child { node {11}
				child { node {5}
						child { node {g,2} }
						child { node {o,3} }
					}
				child { node {k,6} }
			}
		child { node {u,14} };
	\end{tikzpicture}
}

\huffstep{%
	\textbf{5.} Łączymy $19$ i~$e$. \\
	Zostają: $(25), (39)$.
}{%
	\begin{tikzpicture}[htree]
		\node {39}
		child { node {19}
				child { node {m,8} }
				child { node {l,11} }
			}
		child { node {e,20} };
	\end{tikzpicture}
}

\huffstep{%
	\textbf{6.} Łączymy $25$ i~$39$. \\
	(ostateczne drzewo kodów)
}{%
	\begin{tikzpicture}[htree]
		\node {64}
		child { node {25}
				child { node {11}
						child { node {5}
								child { node {g,2} edge from parent node[elabel left] {0} }
								child { node {o,3} edge from parent node[elabel right] {1} }
								edge from parent node[elabel left] {0}
							}
						child { node {k,6} edge from parent node[elabel right] {1} }
						edge from parent node[elabel left] {0}
					}
				child { node {u,14} edge from parent node[elabel right] {1} }
				edge from parent node[elabel left] {0}
			}
		child { node {39}
				child { node {19}
						child { node {m,8} edge from parent node[elabel left] {0} }
						child { node {l,11} edge from parent node[elabel right] {1} }
						edge from parent node[elabel left] {0}
					}
				child { node {e,20} edge from parent node[elabel right] {1} }
				edge from parent node[elabel right] {1}
			};
	\end{tikzpicture}
}

Odczytany kod z~drzewa w~kroku 6:
\begin{flalign*}
	u & \implies 01   &  & \\
	e & \implies 11   &  & \\
	k & \implies 001  &  & \\
	m & \implies 100  &  & \\
	l & \implies 101  &  & \\
	g & \implies 0000 &  & \\
	o & \implies 0001 &  &
\end{flalign*}

\paragraph{Liczba bitów.}
\[
	\boxed{B(T) = 2\cdot 4 + 3\cdot 4 + 6\cdot 3 + 8\cdot 3 + 11\cdot 3 + 14\cdot 2 + 20\cdot 2 = 163 \text{ bitów}}
\]
(wobec $64 \cdot 3 = 192$ bitów przy kodzie stałej długości po $3$ bity na znak).

% ── Rozwiązanie: Szeregowanie zadań z deadlinami -- instancja 2 ────
\subsection{Szeregowanie zadań z deadlinami -- instancja 2}

Dane: $n = 7$ zadań o~jednostkowym czasie wykonania. Kary $w_i$ są już ułożone
nierosnąco ($w_1 \geq \dots \geq w_7$), więc zadania rozważamy w~kolejności
$z_1, \dots, z_7$. Algorytm \textsc{Greedy} dodaje $z_i$ do~$A$ wtedy i~tylko
wtedy, gdy zbiór pozostaje wolny od kar, tzn.\ profil
$N_t(A) = |\{z_j \in A : d_j \leq t\}|$ nie przekracza $(1, 2, \ldots, n)$.
Zadania terminowe wyróżniono na~\textcolor{accentB}{pomarańczowo}.

\begin{aztable}
	\centering
	\renewcommand{\arraystretch}{1.2}
	\begin{NiceTabular}{c|ccccccc}[margin,hvlines]
		$i$   & \color{accentB}1 & 2 & \color{accentB}3 & \color{accentB}4 & \color{accentB}5 & \color{accentB}6 & 7  \\
		\hline
		$d_i$ & \color{accentB}1 & 1 & \color{accentB}3 & \color{accentB}2 & \color{accentB}5 & \color{accentB}4 & 3  \\
		$w_i$ & \color{accentB}60 & 55 & \color{accentB}45 & \color{accentB}35 & \color{accentB}25 & \color{accentB}20 & 10
	\end{NiceTabular}
\end{aztable}

\begin{dygresja}
	Poniżej $A$ wypisane rosnąco po~deadlinach, a~po prawej profil
	$N(A) = (N_1, \ldots, N_7)$, który musi pozostać $\leq (1, 2, 3, 4, 5, 6, 7)$:
	\begin{enumerate}[nosep]
		\item $A = (z_1)$ \hfill $N(A) = (1, 1, 1, 1, 1, 1, 1)$
		\item $z_2$ ($d = 1$) nie dodajemy, bo $N_1(A \cup \{z_2\}) = 2 > 1$.
		\item $A = (z_1, z_3)$ \hfill $N(A) = (1, 1, 2, 2, 2, 2, 2)$
		\item $A = (z_1, z_4, z_3)$ \hfill $N(A) = (1, 2, 3, 3, 3, 3, 3)$
		\item $A = (z_1, z_4, z_3, z_5)$ \hfill $N(A) = (1, 2, 3, 3, 4, 4, 4)$
		\item $A = (z_1, z_4, z_3, z_6, z_5)$ \hfill $N(A) = (1, 2, 3, 4, 5, 5, 5)$
		\item $z_7$ ($d = 3$) nie dodajemy, bo $N_3(A \cup \{z_7\}) = 4 > 3$.
	\end{enumerate}
\end{dygresja}

\paragraph{Harmonogram.}
Zadania terminowe ustawiamy rosnąco po~deadlinach, a~spóźnione ($z_2, z_7$) na
końcu:
\[
	\boxed{(z_1,\ z_4,\ z_3,\ z_6,\ z_5,\ z_2,\ z_7)}
\]
Suma kar zadań spóźnionych: $w_2 + w_7 = 55 + 10 = 65$.

% ── Rozwiązanie: Greedy Set Cover -- instancja 2 ─────────────────
\subsection{Greedy Set Cover -- instancja 2}

Dane: $X = \{1, \ldots, 9\}$ oraz $\mathcal{F} = \{S_1, \ldots, S_5\}$.
Algorytm \textsc{GSC} w~każdej iteracji dokłada do~$\mathcal{C}$ zbiór
maksymalizujący $|S \cap U|$ (remisy na korzyść mniejszego indeksu) i~odejmuje
$U \Assign U \setminus S$. Zbiory wybrane wyróżniono na~\textcolor{accentB}{pomarańczowo}.

\begin{aztable}
	\centering
	\renewcommand{\arraystretch}{1.2}
	\begin{NiceTabular}{c|l|c}[margin,hvlines]
		$S$                  & elementy                & $|S|$ \\
		\hline
		\color{accentB}$S_1$ & \color{accentB}$\{1,2,3,4,5\}$ & \color{accentB}5 \\
		\color{accentB}$S_2$ & \color{accentB}$\{6,7,8,9\}$   & \color{accentB}4 \\
		$S_3$                & $\{1,4,7\}$             & 3 \\
		$S_4$                & $\{2,5,8\}$             & 3 \\
		$S_5$                & $\{3,6,9\}$             & 3
	\end{NiceTabular}
\end{aztable}

\begin{dygresja}
	W~każdym wierszu po~lewej bieżący zbiór niepokrytych $U$, a~po~prawej
	rozważane przecięcia $|S \cap U|$ (pogrubiony zwycięzca):
	\begin{enumerate}[nosep]
		\item $U = \{1,\ldots,9\}$:\quad
		      $|S_1 \cap U| = \mathbf{5}$, $|S_2 \cap U| = 4$,
		      $|S_3 \cap U| = |S_4 \cap U| = |S_5 \cap U| = 3$
		      \hfill wybór $S_1$
		\item $U = \{6,7,8,9\}$:\quad
		      $|S_2 \cap U| = \mathbf{4}$, $|S_5 \cap U| = 2$,
		      $|S_3 \cap U| = |S_4 \cap U| = 1$
		      \hfill wybór $S_2$
	\end{enumerate}
	Po~drugiej iteracji $U = \emptyset$, pętla się kończy.
\end{dygresja}

\paragraph{Wynik.}
Algorytm zachłanny zwraca pokrycie
\[
	\boxed{\mathcal{C} = \{S_1,\ S_2\}}
\]
o~rozmiarze $2$. Ponieważ $S_1 \cup S_2 = X$, a~żaden pojedynczy zbiór nie pokrywa
całego $X$, jest to zarazem pokrycie optymalne ($|\mathcal{C}^*| = 2$) --- tutaj
zachłanność trafia w~optimum.

% ── Rozwiązanie: Approx Vertex Cover -- instancja 2 ──────────────
\subsection{Approx Vertex Cover -- instancja 2}

Algorytm \textsc{AVC} dopóki zostają niepokryte krawędzie, wybiera (tu:
leksykograficznie najmniejszą) krawędź $uv \in F$, dorzuca do~pokrycia $C$
\emph{oba} jej końce i~usuwa z~$F$ wszystkie krawędzie incydentne z~$u$ lub $v$.
Krawędzie wybrane (skojarzenie $B$) na~\textcolor{accentB}{pomarańczowo}, usuwane
wraz z~nimi --- szaro, wierzchołki dokładane do~$C$ --- niebiesko.

% Lokalne style i makra (sufiks B, by nie kolidować z instancją 1).
\tikzset{
	vcnodeB/.style={circle, draw, minimum size=0.55cm, inner sep=1pt, font=\small},
	covB/.style={fill=accent!25},
	selB/.style={accentB, line width=1.2pt},
	goneB/.style={gray!45, dashed},
}
% #1..#7: style wierzchołków a,b,c,d,e,f,g
\newcommand{\vcnodesB}[7]{%
	\node[vcnodeB,#1] (a) at (0,2)    {$a$};
	\node[vcnodeB,#2] (b) at (1.5,2)  {$b$};
	\node[vcnodeB,#3] (c) at (3,2)    {$c$};
	\node[vcnodeB,#4] (d) at (4.5,2)  {$d$};
	\node[vcnodeB,#5] (e) at (0.75,0) {$e$};
	\node[vcnodeB,#6] (f) at (2.25,0) {$f$};
	\node[vcnodeB,#7] (g) at (3.75,0) {$g$};
}
% Makro może mieć najwyżej 9 parametrów (#10 = #1 + "0"), więc dzielimy na dwa.
% \vcedgesBa: ab,bc,cd,ae,be   \vcedgesBb: bf,cf,cg,dg,ef
\newcommand{\vcedgesBa}[5]{%
	\draw[#1] (a)--(b); \draw[#2] (b)--(c); \draw[#3] (c)--(d); \draw[#4] (a)--(e);
	\draw[#5] (b)--(e);}
\newcommand{\vcedgesBb}[5]{%
	\draw[#1] (b)--(f); \draw[#2] (c)--(f); \draw[#3] (c)--(g); \draw[#4] (d)--(g);
	\draw[#5] (e)--(f);}
\newcommand{\vcpanelB}[1]{%
	\begin{tikzpicture}[scale=0.55, baseline=(current bounding box.center)]#1\end{tikzpicture}}
\newcommand{\vclabelB}[1]{\node[font=\footnotesize] at (2.25,-0.9) {#1};}

\begin{dygresja}
	Przebieg pętli ($F$ -- niepokryte krawędzie, $B$ -- wybierane krawędzie):
	\begin{enumerate}[nosep]
		\item $F = \{ab, ae, bc, be, bf, cd, cf, cg, dg, ef\}$. Wybór $ab$; usuwamy
		      $ab, ae, bc, be, bf$ (incydentne z~$a$ lub $b$).
		\item $F = \{cd, cf, cg, dg, ef\}$. Wybór $cd$; usuwamy $cd, cf, cg, dg$
		      (incydentne z~$c$ lub $d$).
		\item $F = \{ef\}$. Wybór $ef$; usuwamy $ef$ --- $F = \emptyset$.
	\end{enumerate}
\end{dygresja}

\begin{azfigure}
	\centering
	\vcpanelB{\vcnodesB{}{}{}{}{}{}{}\vcedgesBa{}{}{}{}{}\vcedgesBb{}{}{}{}{}\vclabelB{start}}%
	\;$\rightarrow$\;%
	\vcpanelB{\vcnodesB{covB}{covB}{}{}{}{}{}\vcedgesBa{selB}{goneB}{}{goneB}{goneB}\vcedgesBb{goneB}{}{}{}{}\vclabelB{wybór $ab$}}%
	\;$\rightarrow$\;%
	\vcpanelB{\vcnodesB{covB}{covB}{covB}{covB}{}{}{}\vcedgesBa{selB}{goneB}{selB}{goneB}{goneB}\vcedgesBb{goneB}{goneB}{goneB}{goneB}{}\vclabelB{wybór $cd$}}%
	\;$\rightarrow$\;%
	\vcpanelB{\vcnodesB{covB}{covB}{covB}{covB}{covB}{covB}{}\vcedgesBa{selB}{goneB}{selB}{goneB}{goneB}\vcedgesBb{goneB}{goneB}{goneB}{goneB}{selB}\vclabelB{wybór $ef$}}
\end{azfigure}

\paragraph{Wynik.}
Algorytm dokłada oba końce trzech wybranych krawędzi $B = \{ab, cd, ef\}$,
zwracając pokrycie
\[
	\boxed{C = \{a,\ b,\ c,\ d,\ e,\ f\}}
\]
o~liczności $|C| = 2|B| = 6$ (wierzchołek $g$ nie wszedł). Optymalne pokrycie ma
rozmiar $4$, np.\ $C^* = \{b, c, d, e\}$: każda krawędź ma koniec w~$b$, $c$, $d$
lub~$e$. Iloraz $|C| / |C^*| = 6/4 = 1{,}5 \leq 2$, zgodnie z~gwarancją
aproksymacji \textsc{AVC}.

% ── Rozwiązanie: Exact Subset Sum -- instancja 2 ─────────────────
\subsection{Exact Subset Sum -- instancja 2}

Dane: $S = \{2, 3, 7, 8, 10\}$ (w~kolejności $x_1, \ldots, x_5$) oraz $t = 15$.
Algorytm \textsc{ExactSubsetSum} startuje z~$L_0 = (0)$ i~w~$i$-tej iteracji scala
$L_{i-1}$ z~$L_{i-1} + x_i$, po~czym usuwa elementy $> t$. Element wprowadzony
przez przesunięcie $L_{i-1} + x_i$ wyróżniono na~\textcolor{accentB}{pomarańczowo}.

\begin{dygresja}
	Przebieg iteracji ($L_{i-1} + x_i$ to składnik przesunięty):
	\begin{enumerate}[nosep]
		\item $x_1 = 2$:\quad
		      $(0) \cup (\textcolor{accentB}{2}) = (0, \textcolor{accentB}{2})$
		      \hfill $L_1 = (0, 2)$
		\item $x_2 = 3$:\quad
		      $(0, 2) \cup (\textcolor{accentB}{3}, \textcolor{accentB}{5})
		      = (0, 2, \textcolor{accentB}{3}, \textcolor{accentB}{5})$
		      \hfill $L_2 = (0, 2, 3, 5)$
		\item $x_3 = 7$:\quad scalenie z~$(7, 9, 10, 12)$, nic $> 15$
		      \hfill\\
		      $L_3 = (0, 2, 3, 5, \textcolor{accentB}{7}, \textcolor{accentB}{9},
		      \textcolor{accentB}{10}, \textcolor{accentB}{12})$
		\item $x_4 = 8$:\quad przesunięcie $L_3 + 8$ daje
		      $8, 10, 11, 13, 15, 17, 18, 20$; po~usunięciu $> 15$ dochodzą
		      $\textcolor{accentB}{8}, \textcolor{accentB}{11}, \textcolor{accentB}{13},
		      \textcolor{accentB}{15}$ ($10$ już było) \\
		      $L_4 = (0, 2, 3, 5, 7, \textcolor{accentB}{8}, 9, 10,
		      \textcolor{accentB}{11}, 12, \textcolor{accentB}{13}, \textcolor{accentB}{15})$
		\item $x_5 = 10$:\quad przesunięcie $L_4 + 10$ daje
		      $10, 12, 13, 15, 17, \ldots$; wszystkie $\leq 15$ ($10, 12, 13, 15$) już
		      są na liście, więc \emph{nic nowego}
		      \hfill $L_5 = L_4$
	\end{enumerate}
\end{dygresja}

\paragraph{Wynik.}
Największy element listy $L_5$ nieprzekraczający $t = 15$ to
\[
	\boxed{z = 15}
\]
czyli dokładnie $t$ -- istnieje podzbiór o~sumie równej progowi, np.\
$\{7, 8\}$ lub $\{2, 3, 10\}$. Wersja decyzyjna ma więc odpowiedź \emph{tak},
a~optymalna suma $\leq t$ wynosi $15$.

% ── Rozwiązanie: Approx Subset Sum -- instancja 2 ────────────────
\subsection{Approx Subset Sum -- instancja 2}

Dane: $S = \{10, 12, 17, 25\}$ (w~kolejności $x_1, \ldots, x_4$), $t = 50$ oraz
$\varepsilon = 0{,}5$. Wtedy $n = 4$ i~próg przycinania
$\delta = \frac{\varepsilon}{2n} = \frac{0{,}5}{8} = 0{,}0625$ (mnożnik
$1 + \delta = 1{,}0625$). Startujemy z~$L_0 = (0)$.

W~każdej iteracji najpierw scalamy $L_{i-1}$ z~$L_{i-1} + x_i$, potem usuwamy
elementy $> t = 50$, a~na~końcu przycinamy: zachowujemy $y$, gdy
$y > last \cdot 1{,}0625$, gdzie $last$ to ostatnio zachowany element.
\begin{dygresja} \leavevmode
	\begin{itemize}[nosep]
		\item \textbf{$i = 1$, $x_1 = 10$.} \\
		      scalenie: $(0) \cup (10) = (0, 10)$; po~usunięciu $> 50$: bez zmian; \\
		      przycinanie: \\
		      \begin{tblr}{colspec={r c l l}, column{4}={font=\itshape}, rowsep=1pt}
			      $0$  &     &     & zachowane, $last = 0$  \\
			      $10$ & $>$ & $0$ & zachowane, $last = 10$ \\
		      \end{tblr} \\
		      $L_1 = (0, 10)$.
		\item \textbf{$i = 2$, $x_2 = 12$.} \\
		      scalenie: $(0, 10) \cup (12, 22) = (0, 10, 12, 22)$;
		      po~usunięciu $> 50$: bez zmian; \\
		      przycinanie: \\
		      \begin{tblr}{colspec={r c l l}, column{4}={font=\itshape}, rowsep=1pt}
			      $0$  &     &                                & zachowane, $last = 0$  \\
			      $10$ & $>$ & $0$                            & zachowane, $last = 10$ \\
			      $12$ & $>$ & $10 \cdot 1{,}0625 = 10{,}625$ & zachowane, $last = 12$ \\
			      $22$ & $>$ & $12 \cdot 1{,}0625 = 12{,}75$  & zachowane, $last = 22$ \\
		      \end{tblr} \\
		      $L_2 = (0, 10, 12, 22)$.
		\item \textbf{$i = 3$, $x_3 = 17$.} \\
		      scalenie: $(0, 10, 12, 22) \cup (17, 27, 29, 39)
			      = (0, 10, 12, 17, 22, 27, 29, 39)$;
		      po~usunięciu $> 50$: bez zmian; \\
		      przycinanie: \\
		      \begin{tblr}{colspec={r c l l}, column{4}={font=\itshape}, rowsep=1pt}
			      $0$  &     &                                 & zachowane, $last = 0$  \\
			      $10$ & $>$ & $0$                             & zachowane, $last = 10$ \\
			      $12$ & $>$ & $10 \cdot 1{,}0625 = 10{,}625$  & zachowane, $last = 12$ \\
			      $17$ & $>$ & $12 \cdot 1{,}0625 = 12{,}75$   & zachowane, $last = 17$ \\
			      $22$ & $>$ & $17 \cdot 1{,}0625 = 18{,}0625$ & zachowane, $last = 22$ \\
			      $27$ & $>$ & $22 \cdot 1{,}0625 = 23{,}375$  & zachowane, $last = 27$ \\
			      $29$ & $>$ & $27 \cdot 1{,}0625 = 28{,}6875$ & zachowane, $last = 29$ \\
			      $39$ & $>$ & $29 \cdot 1{,}0625 = 30{,}8125$ & zachowane, $last = 39$ \\
		      \end{tblr} \\
		      $L_3 = (0, 10, 12, 17, 22, 27, 29, 39)$.
		\item \textbf{$i = 4$, $x_4 = 25$.} \\
		      scalenie \\
		      $(0, 10, 12, 17, 22, 27, 29, 39) \cup (25, 35, 37, 42, 47, 52, 54, 64)$ \\
		      $ = (0, 10, 12, 17, 22, 25, 27, 29, 35, 37, 39, 42, 47, 52, 54, 64)$; \\
		      po~usunięciu $> 50$ odpadają $52, 54, 64$:
		      $(0, 10, 12, 17, 22, 25, 27, 29, 35, 37, 39, 42, 47)$; \\
		      przycinanie: \\
		      \begin{tblr}{colspec={r c l l}, column{4}={font=\itshape}, rowsep=1pt}
			      $0$  &        &                                 & zachowane, $last = 0$  \\
			      $10$ & $>$    & $0$                             & zachowane, $last = 10$ \\
			      $12$ & $>$    & $10{,}625$                      & zachowane, $last = 12$ \\
			      $17$ & $>$    & $12{,}75$                       & zachowane, $last = 17$ \\
			      $22$ & $>$    & $17 \cdot 1{,}0625 = 18{,}0625$ & zachowane, $last = 22$ \\
			      $25$ & $>$    & $22 \cdot 1{,}0625 = 23{,}375$  & zachowane, $last = 25$ \\
			      $27$ & $>$    & $25 \cdot 1{,}0625 = 26{,}5625$ & zachowane, $last = 27$ \\
			      $29$ & $>$    & $27 \cdot 1{,}0625 = 28{,}6875$ & zachowane, $last = 29$ \\
			      $35$ & $>$    & $29 \cdot 1{,}0625 = 30{,}8125$ & zachowane, $last = 35$ \\
			      $37$ & $\leq$ & $35 \cdot 1{,}0625 = 37{,}1875$ & usunięte               \\
			      $39$ & $>$    & $37{,}1875$                     & zachowane, $last = 39$ \\
			      $42$ & $>$    & $39 \cdot 1{,}0625 = 41{,}4375$ & zachowane, $last = 42$ \\
			      $47$ & $>$    & $42 \cdot 1{,}0625 = 44{,}625$  & zachowane, $last = 47$ \\
		      \end{tblr} \\
		      $L_4 = (0, 10, 12, 17, 22, 25, 27, 29, 35, 39, 42, 47)$.
	\end{itemize}
\end{dygresja}

\paragraph{Wynik.}
Zwracamy największy element listy $L_4$:
\[
	\boxed{z = 47}
\]
realizowany podzbiorem $\{10, 12, 25\}$. Tu optimum dokładne to również
$z^* = 47$ (ta sama suma), więc $\frac{z^*}{z} = 1 \leq 1 + \varepsilon = 1{,}5$
--- przycięcie usunęło jedynie $37$, nie naruszając optimum.

% ── Rozwiązanie: Maximum Matching -- instancja 2 ─────────────────
\subsection{Maximum Matching -- instancja 2}

Graf dwudzielny $G = (V_1, V_2; E)$ z~$V_1 = \{a, b, c, d, e\}$,
$V_2 = \{p, q, r, s, t\}$ oraz krawędziami $ap, aq, bp, cq, cr, dr, ds, dt, es$.
Algorytm \textsc{MM} startuje z~$M = \emptyset$ i~dopóki \textsc{APS} zwraca
ścieżkę powiększającą $P$, podmienia $M \Assign M \div E(P)$. \textsc{APS}
przeszukuje $G_M$ w~głąb (sąsiedzi alfabetycznie) z~wolnego wierzchołka $V_1$
do~wolnego w~$V_2$. Na~rysunkach pierścień otacza wierzchołki \emph{wolne},
a~strzałki pokazują orientację krawędzi w~$G_M$ (skojarzone w~górę, do~$V_1$;
wolne w~dół, do~$V_2$). Każdy wiersz pokazuje $G_M$ ze~znalezioną ścieżką
powiększającą (\textcolor{path}{na~różowo}), a~po~strzałce $\Rightarrow$ ---
wynikowe skojarzenie $M$ (\textcolor{accentB}{na~pomarańczowo}).

% Lokalne makra (sufiks B, by nie kolidować z instancją 1).
\tikzset{
	mmnodeB/.style={circle, fill=black, inner sep=1.5pt},
	selB/.style={accentB, line width=1.4pt, <-}, % skojarzona: strzałka V_2->V_1 (w górę)
	augFB/.style={path, line width=1.4pt},       % na ścieżce, wolna (w dół, domyślne ->)
	augMB/.style={path, line width=1.4pt, <-},   % na ścieżce, skojarzona (w górę)
}
% Strzałki w G_M: krawędź wolna ->, skojarzona przejmuje <- ze stylu selB.
% Pierścień wokół wierzchołków wolnych (lista w #1, pusta dozwolona).
\newcommand{\mmfreeRB}[1]{\ifx\relax#1\relax\else
	\foreach \n in {#1} {\draw[black] (\n) circle (3.5pt);}\fi}
% #1..#9: style krawędzi ap,aq,bp,cq,cr,dr,ds,dt,es (sel albo puste)
\newcommand{\mmedgesB}[9]{%
	\foreach \x/\n in {0/a, 1/b, 2/c, 3/d, 4/e} {\node[mmnodeB, label=above:$\n$] (\n) at (\x,1) {};}
	\foreach \x/\n in {0/p, 1/q, 2/r, 3/s, 4/t} {\node[mmnodeB, label=below:$\n$] (\n) at (\x,0) {};}
	\draw[->,#1] (a)--(p); \draw[->,#2] (a)--(q); \draw[->,#3] (b)--(p); \draw[->,#4] (c)--(q);
	\draw[->,#5] (c)--(r); \draw[->,#6] (d)--(r); \draw[->,#7] (d)--(s); \draw[->,#8] (d)--(t);
	\draw[->,#9] (e)--(s);}
% #1: style krawędzi, #2: etykieta, #3: lista wolnych, #4: początek ścieżki (pierścień różowy, pusty dozwolony)
\newcommand{\mmpanelB}[4]{%
	\begin{tikzpicture}[x=0.7cm, y=0.85cm, scale=0.72, >=stealth, baseline=(current bounding box.center)]
		\mmedgesB#1\mmfreeRB{#3}%
		\ifx\relax#4\relax\else\draw[path, line width=1.4pt] (#4) circle (3.5pt);\fi
		\node[font=\footnotesize] at (2,-0.95) {#2};\end{tikzpicture}}

\begin{dygresja}
	Przebieg pętli (ścieżki powiększające na~\textcolor{path}{różowo},
	krawędzie w~skojarzeniu na~\textcolor{accentB}{pomarańczowo}):
	\begin{enumerate}[nosep]
		\item $M = \emptyset$. \textsc{APS} z~$a$ znajduje wolną krawędź $P_1 = a\,p$.
		      \hfill $M = \{ap\}$
		\item Z~$b$: wolna $b\!\to\!p$, skojarzona $p\!\to\!a$, wolna $a\!\to\!q$
		      ($q$ wolne). Ścieżka $P_2 = b\,p\,a\,q$ ($ap$ wychodzi, $bp, aq$ wchodzą).
		      \hfill $M = \{aq, bp\}$
		\item Z~$c$: $c\!\to\!q$ prowadzi do martwego punktu, więc bierzemy
		      $c\!\to\!r$ ($r$ wolne): $P_3 = c\,r$.
		      \hfill $M = \{aq, bp, cr\}$
		\item Z~$d$: $d\!\to\!r$ prowadzi do martwego punktu, więc $d\!\to\!s$
		      ($s$ wolne): $P_4 = d\,s$.
		      \hfill $M = \{aq, bp, cr, ds\}$
		\item Z~$e$: $e\!\to\!s$, skojarzona $s\!\to\!d$, wolna $d\!\to\!t$ ($t$
		      wolne). Ścieżka $P_5 = e\,s\,d\,t$ ($ds$ wychodzi, $es, dt$ wchodzą).
		      \hfill $M = \{aq, bp, cr, es, dt\}$
	\end{enumerate}
	Brak wolnych wierzchołków w~$V_1$ -- \textsc{APS} zwraca NULL, pętla kończy się.
\end{dygresja}

% Każdy wiersz: G_M ze ścieżką powiększającą (\textcolor{path}{na różowo}, pierścień
% startowy też różowy) -> wynikowe M. Kolumny w tblr wyrównują rysunki w pionie.
\begin{azfigure}
	\centering
	\begin{tblr}{colspec={ccc}, rowsep=10pt, colsep=6pt}
		\mmpanelB{{augFB}{}{}{}{}{}{}{}{}}{$P_1 = a\,p$}{a,b,c,d,e,p,q,r,s,t}{a}
		& $\Rightarrow$ &
		\mmpanelB{{selB}{}{}{}{}{}{}{}{}}{$M = \{ap\}$}{b,c,d,e,q,r,s,t}{} \\
		\mmpanelB{{augMB}{augFB}{augFB}{}{}{}{}{}{}}{$P_2 = b\,p\,a\,q$}{b,c,d,e,q,r,s,t}{b}
		& $\Rightarrow$ &
		\mmpanelB{{}{selB}{selB}{}{}{}{}{}{}}{$M = \{aq, bp\}$}{c,d,e,r,s,t}{} \\
		\mmpanelB{{}{selB}{selB}{}{augFB}{}{}{}{}}{$P_3 = c\,r$}{c,d,e,r,s,t}{c}
		& $\Rightarrow$ &
		\mmpanelB{{}{selB}{selB}{}{selB}{}{}{}{}}{$M = \{aq, bp, cr\}$}{d,e,s,t}{} \\
		\mmpanelB{{}{selB}{selB}{}{selB}{}{augFB}{}{}}{$P_4 = d\,s$}{d,e,s,t}{d}
		& $\Rightarrow$ &
		\mmpanelB{{}{selB}{selB}{}{selB}{}{selB}{}{}}{$M = \{aq, bp, cr, ds\}$}{e,t}{} \\
		\mmpanelB{{}{selB}{selB}{}{selB}{}{augMB}{augFB}{augFB}}{$P_5 = e\,s\,d\,t$}{e,t}{e}
		& $\Rightarrow$ &
		\mmpanelB{{}{selB}{selB}{}{selB}{}{}{selB}{selB}}{$M = \{aq, bp, cr, dt, es\}$}{}{} \\
	\end{tblr}
\end{azfigure}

\paragraph{Wynik.}
Algorytm zwraca skojarzenie
\[
	\boxed{M = \{aq,\ bp,\ cr,\ dt,\ es\}}
\]
o~liczności $|M| = 5$. Jest ono \emph{doskonałe} --- każdy z~$5$ wierzchołków obu
stron jest skojarzony --- więc i~najliczniejsze. Ścieżki powiększające $P_2$
i~$P_5$ pokazują istotę metody: przekładają skojarzenie wzdłuż naprzemiennych
krawędzi i~zwiększają $|M|$ o~$1$, odblokowując kolejne wierzchołki.


\end{document}
